% 原书第 491 页（PDF 第 504 页）
\BookPage{491}{504}
\BookHeading{何震年表}{何震年表}{万仕国*}
\noindent 1886年一岁

何震出生于扬州。

据刘师培〈清故武进县学教谕何先生墓志铭〉称:「其先出自东海，奕业载德，光于方志。」何氏世居扬州，占籍仪征，遂称仪征人。

曾祖父口口，县学生员。

祖父何域，诰赠知府。

父何承霖（1838-1897），字春海，兄弟八人，行五。以县学生员中丁卯科举人，补授八旗官学教习。刘师培〈清故武进县学教谕何先生墓志铭〉称「会以大挑，再授知县。」「补常州武进县学教谕。」光绪二十四年七月卒于仪征十二圩舟次，「春秋六十」。有女子四人，二男、二女。

长女何采苹（1869-1916），即所谓「何大姑」，刘师培〈何大姑哀赞〉称其「伤足弱行，因贞不字」，随父「从官武进一十二年，以孝烝烝，型其群季」。「春秋四十有八，民国五年口月，卒于弟家辂之居」。

长子何家辂，字质卿，号誉生，增生，娶刘师培伯姊、刘谦甫之女刘师韫（字谢仪）。上海图书馆藏《仪征县学生名册》抄本，于「光绪二十三年丁酉科」有其入县学的名字。善书法，刘富曾〈亡侄师培墓志铭〉即为其书丹，另有书法作品传世。其余无考。

\vfill
\begin{BookNotes}
\BookNote{*}{作者任职于江苏省仪征市人大常委会办公室。}
\end{BookNotes}

% 原书第 492 页（PDF 第 505 页）
\BookPage{492}{505}
次子何家勤，生平亦无可考。

季女何震，原名班。〈何大姑哀赞〉称:「震年及龀，授以工言。阙则弥缝，灾则匡救。」

1897年十二岁

父何承霖卒。何采苹回到扬州，奉养母亲。〈何大姑哀赞〉称:「会遘父忧，丧存甯戚，言归奉母，好是幽闲。德慧在躬，不违终日。」

刘师培〈何大姑哀赞〉以何震的口气称:「我丧严亲，二六之年。」梅〈青溪旧屋仪征刘氏五世小记〉称，何班「幼年在家，秉承闺训甚严，不见生人」。

1904年十九岁

6月，何班与刘师培在扬州结婚。因为刘师培的堂姊刘师韫已经嫁给了何家辂，所以这门亲上加亲的婚事，在扬州被称作「还门亲」，是扬州旧时代大家族之间常见的一种婚姻现象。在刘师培的宣传影响下，何班「结婚后思想大为解放，以后就与舅氏（指刘师培）每出必同行了」1。

7月，何班随刘师培赴上海，改名何震，字志剑，入爱国女学，并与林宗素(1878-1944)等新派女性相交往。7月26日，何震在《警钟日报》发表〈赠侯官林宗素女士〉诗，称：「献身甘作苏菲亚，爱国群推玛利侬。言念神州诸女杰，何时杯酒饮黄龙。」林宗素在此诗后附识称扬道:「何女士为刘申叔先生夫人，结婚才逾月。先生于吾国学界为有数之人物，其夫人学问宗旨足以称之。吾为吾国女界贺！吾为刘先生贺！」可见刘师培夫妇在当时进步人士看来，是非常相配的一对。而且，从林宗素的褒扬之辞看来，何震至少是通文墨的，否则林氏所恭维的「学问宗旨足以称之」，就成了对何震的嘲笑之言了。

11月19日（农历十月二十五），刘师培与万福华行刺王之春于上海英租界四马

\vfill
\begin{BookNotes}
\BookNote{1.}{梅着、梅英超整理，《青溪旧屋仪征刘氏五世小记》，上海古籍出版社，2004年，页36。}
\end{BookNotes}

% 原书第 493 页（PDF 第 506 页）
\BookPage{493}{506}
\noindent 路、湖北路口的金谷香番菜馆，功败垂成，刘师培等被拘，旋出狱，并与林獬等集资延聘律师，为万福华辩护。万福华后被判刑十年。

1905年二十岁

1月，刘师培与邓实、黄节等在上海发起成立国学保存会，并编辑出版《国粹学报》，「以发明国学、保存国粹为宗旨」。

3月，清政府查封《警钟日报》，刘师培等被通缉，乃改名金少甫，逃亡嘉兴，藏匿敖嘉熊家，并助理温台处会馆事宜。2何震则返回扬州避难。

在上海期间，何震究竟对刘师培所从事的推翻满清斗争中起了什么样的作用，现在没有任何资料可以说明。但从刘师培，参与行刺王之春、加入同盟会等经历看，何震至少没有成为阻力，这应当是可以肯定的。

1906年二十一岁

春，刘师培应陈独秀等招，赴安徽芜湖，任安徽公学、皖江中学教师。何震与婆母李汝同往，住芜湖弋矶山。

刘师培在芜湖期间，援引敖嘉熊、徐慕达、苏曼殊等赴芜湖，与陈独秀、陶成章、邓绳侯、龚宝铨等继续宣传排满革命，参与组织暗杀团体黄氏学校，积极发展同盟会员，并托胡渭清将一百名新吸收的同盟会会员名册送交上海蔡元培。刘师培等在芜湖的革命活动引起清廷注意，「端方曾告知蒯光典，他打算拿办安徽公学革命党人，经蒯婉言解释，端方才没有下令实行」3。

1907年二十二岁

2月13日（农历元旦），因国内对革命党人迫害日烈，刘师培拟在上海创办国

\vfill
\begin{BookNotes}
\BookNote{2.}{冯自由，〈兴中会时期之革命同志〉，参见氏着《革命逸史》（中），新星出版社，2009年。}
\BookNote{3.}{〈辛亥前安徽文教界的革命活动〉，全国政协文史资料研究委员会编《辛亥革命回忆录》（四)，中华书局，1963年。}
\end{BookNotes}

% 原书第 494 页（PDF 第 507 页）
\BookPage{494}{507}
\noindent 粹学堂的计划因资金无着而流产，于是听从马君武建议，应章太炎之邀，与何震、汪公权、苏曼殊同赴日本。到日本后，何震开始公开参与社会活动，业余时间从苏曼殊学绘画，称女弟子。

3月，民报社发生「倒孙风潮」，刘师培与孙中山等同盟会本部成员产生隔闵。

4月，何震与刘师培、章太炎、张继、苏曼殊、陈独秀等共同发起成立「亚洲和亲会」。该会以「反对帝国主义，期使亚洲已失主权之民族各得独立」为宗旨，「除主张侵略主义者，无论民族主义、共和主义、无政府主义，皆得入会」，广泛团结亚洲各国志士，「互相扶持，使各得独立」。「凡亚洲诸国，若一国有革命事，余国同会者应互相协助，不论直接、间接，总以功能所及为限。」4陶铸〈亚洲和亲会约章中文抄稿附识〉云:「此会成立于1907（光绪三十三年）丁未之春，首由中印两国革命志士发起于日本之东京。《亚洲和亲会约章》为章太炎先生之手笔，译成英文。开章明义，即为反对帝国主义。其后陆续加入者有：越南、缅甸、菲律宾、朝鲜诸邦，形成亚洲民族解放统一战线，以余记忆所及，中国方面入会者有：章太炎（炳麟）、张溥泉（继）、刘申叔（师培）、何殷振（震）、苏子谷（元瑛，法名曼殊）、陈仲甫（独秀）、吕剑秋（复）、罗黑子（象陶）及余等数十人。」

夏，何震在与日本社会党人的接触中，开始接受盛行于日本的无政府主义和社会主义思想的影响，醉心于斯宾塞《女权篇》，提倡男女平等同权、父母姓氏并重，遂改名何殷震。冯自由〈记刘光汉变节始末〉云:「光汉夫妇时与章太炎、苏曼殊同居，渐与日本社会党人之急激派北辉次郎、和田三郎等游，遂醉心社会主义。寻发刊一《天义报》月刊，极力鼓吹社会主义学说，是为我国人发刊社会主义机关报之嚆矢。其妻何震更提倡父母姓并重之说，自号其姓名曰何殷震。」

\vfill
\begin{BookNotes}
\BookNote{4.}{章太炎，〈亚洲和亲会约章〉，引自汤志钧着《章太炎年谱长编》（上)，中华书局，1977年。}
\BookNote{5.}{冯自由，〈革命逸史〉（上)，新星出版社，2009年，页332。}
\end{BookNotes}

% 原书第 495 页（PDF 第 508 页）
\BookPage{495}{508}
6月，何震与陆恢权、徐亚尊、周怒涛、张旭等发起创办「女子复权会」，由何震出面，创办《天义》作为机关刊物，社址在东京小石川区久坚町二十七番瑜伽师地，以何震的住宅东京牛込区新小川町二丁目八番地作为联络机关，热心宣传女权主义和无政府主义理论。《复报》第十期刊发〈天义报启〉云:「地球之上，邦国环立，然自有人类以来，无一事合于真公。异族之欺陵，君民之悬隔，贫富之差殊，此咸事之属于不公者也。自民族主义明，然后受制于异族者，人人均以为辱；自民约之论昌，然后受制于暴君者，人人均引为耻；自社会主义明，然后受制于富民者，人人均以为羞。由是种族革命、政治革命、经济革命，遂为人民天赋之权。然环顾世界各邦，其实行种族革命者尚占多数。若政治一端，虽实行共和政治者，犹不能尽人而平等；经济一端，更无论矣。试推其原因，则以世界故有之社会，均属于阶级制度，合无量不公不平之习惯，相积而成。故无论其迁变之若何，均含有不平之性质。非破坏固有之社会，决不能扫除阶级，使之尽合于公。顾今之论者，所言之革命，仅以经济革命为止。不知世界固有之阶级，以男女阶级为严。无论东洋有尊男轻妇之风也，即西洋各国号为男女平等者，然服官议政之权，均为女子所无，则是女子所有之权，并贱民而不若。更反观之于中国，则夫可多妻，妻不可多夫；男可再娶，女不可再嫁；服丧则一斩一期，宾祭则此先彼后。即有号为均平者，既嫁之后，内夫家而外母家，所生子女，用父姓而遗母姓，又安得谓之公平乎？夫男女之间，其制度失平且若此，于此而欲破坏固有之阶级，不亦难乎！故欲破社会固有之阶级，必自破男女阶级始。所谓破男女阶级者，即无论男女，均与以相当之教养、相当之权利，使女子不致下于男，男子不能加于女，男对于女若何，即女对于男亦若何。如有女下男而男加女者，则女界共起而诛之，务使相平而后已。夫以男女阶级之严，行之数千载，今也一旦而破之，则凡破坏社会之方法，均可顺次而施行，天下岂有不破之阶级哉！夫居今日之世界，非尽破固有之阶级，不得使之反于公；居今日之中国，非男女革命与种族、政治、经济诸革命并行，亦不得合于真公。震等目击心伤，故创为女子复权会，讨论斯旨，以冀实行其目的。又虑此理之不能共喻也，故刊《天义旬报》，以作本会之机关。惟经营伊始，财政拮

\vfill
\begin{BookNotes}
\BookNote{6.}{陆灵素，名守民，字恢权，又号繁霜，南社女诗人，南社社员、上海华泾的刘季平（江南刘三）之妻，能诗善歌，长于南北曲。刘季平死后，陆灵素搜集整理其遗稿，1946年刊印出油印本，马叙伦为刘三作传。1996年，正式出版《黄叶楼遗稿》一书。}
\end{BookNotes}

% 原书第 496 页（PDF 第 509 页）
\BookPage{496}{509}
\noindent 据，世有赞成此旨者，尚祈慨解囊金，共襄此举，使公平之真理得以普及于寰区，此则世界之幸也。」并说明该报宗旨是:「以破坏固有之社会，实行人类之平等为宗旨，于提倡女界革命外，兼提倡种族、政治、经济诸革命。」

6月10日，《天义报》半月刊第一号出版。

其「图画」栏发表「震敬绘」的女娲像，并题有〈女娲像赞〉云:「于穆娲皇，厥姓惟风，断鼇足，杀黑龙。先禹有功抑下鸿，辟除民害逐共工。是宜报功崇德，与轩羲并隆。」

「社说」栏载署名「何殷震」的〈女子宣布书〉，称:「呜呼！世界之男女，其不平等也久矣。印度之女，自焚以殉男，日本之女，卑屈以事男。欧美各国，虽行一夫一妻之制，号为平等，然议政之权，选举之权，女子均鲜得干预。所谓平权者，果安在邪？更反观之吾中国，则男子之视女子也，几不以人类相待。」何震把古代妇女地位低下的形成，归结为两个历史阶段。「上古之民，战胜他族，则系累其女，械系其身，以为妃妾，由是男为主而女为奴。是为剽掠妇女之时代。继因剽劫易起争端，乃创为俪皮之礼。故古礼所言纳探、纳征，均沿财昏之俗，盖视女子为财产之一也。由是男为人而女为物，是为买卖妇女之时代。积此二因，由是男女之间，遂不平等。」

何震认为，古代男女不平等，至少表现在四个方面。「一曰嫁娶上之不平等。古代之时，位愈尊者妻愈众。…而后世之嫔妃，则更无限制。贵显之家，蓄妾尤众。」「二曰名分上之不平等。男权既伸，其防范女子亦日严，创『一与之齐、终身不改』之说，使女子终事一夫；又谓夫尊妻卑，夫犹天而妻犹地，妻不去夫，犹地不得去天。由是爵则从夫，姓则从夫，而谥亦从夫:以女子为男子附属物。「三曰职务上之不平等。中国『妇』训为『服』，象持帚之形。而《礼记·曲礼篇》亦言:『纳女子于诸侯曰备酒浆，于大夫曰备洒扫。』是古代之妇人，仅以服从为义务。又创为女子不逾阈之说，以禁其自由。后世以降，为女子者，舍治家而外无职务，以有才为大戒，以卑屈为当然。」「四曰礼制上之不平等。夫之于妻，仅服期丧；而妻之于夫，则服丧三年。非惟为夫服重丧也，即夫之父母，亦为之服斩衰；于己之父母，转降为齐衰，非所谓厚于所薄、薄于所厚者邪？且古代之时，


% 原书第 497 页（PDF 第 510 页）
\BookPage{497}{510}
\noindent 父存母殁，为母服齐衰，尤为失理之尤。」

令何震感到特别痛心的是:「夫以男陵女，犹可言也。女子而甘于自屈，抑独何心？岂非社会之习惯、腐儒之学术，有以籍制之邪？」因此她向女界同胞呼吁：「男子者，女子之大敌也。女子一日不与男子平等，则此恨终不磨。」何震认为「女界所应争」之权利，至少包括七个方面:一曰实行一夫一妻之制；二曰既嫁之后，不从夫姓；三曰为父母者，俱男女并重；四曰男女自初生以后，即与以相等之养育；五曰如夫妇既昏而不谐，则告分离；六曰以初昏之男配初昏之女；七曰废尽天下之娼寮，去尽天下之娼女，以扫荡淫风。

何震认为:「以上七事，非女子欲争权利也，特以天赋之权，男女所同；男女同为人类，若不能平等，是为不公，是为背天理，故子女之所争，仅以至公为止境」

针对世人之所疑者三事，何震逐一作了阐述:「一曰女子有生育之苦而生子以还，又有鞠养之劳，故职务不可与男同。不知吾所倡者，非仅女界革命，乃社会革命也。特以女界革命为社会革命之一端。社会革命既实行，所生女子既生以后，即入公设育婴所，不必自为养育也。无养育子女之劳，所尽职务，自可与男相等。二曰女多而男少，则行一夫一妻之制，转属不公。不知女子之多，由于女子不从事战争，而战争之役悉属于男。男子日有死亡，故女多男少。今为女子者，与其甘心为妾受辱而死，曾不若实行破坏，死于沙场，犹有日后之荣名也。故女子果实行革命，事平以后，女子之人数，必与男子之数相等。三曰男既多妻，女亦可多夫，以相抵制。不知女界欲求平等，非徒用抵制之策已也。必以暴力强制男子，使彼不得不与己平。且男子多妻，男子之大失也。今女子亦举而效之，何以塞男子之口乎？况女子多夫，莫若娼妓。今倡多夫之说者，名为抵制男子，实则便其私欲，以蹈娼妓之所为，此则女界之贼也。」

何震强调:「男女同为人类，凡所谓男性、女性者，均习惯使然，教育使然。若不于男女生异视之心，鞠养相同，教育相同，则男女所尽职务亦必可以相同，而男性、女性之名词，直可废灭。此诚所谓男女平等也。近日中国之女子欲争此境，凡种族、政治、经济诸革命，均宜先男子着鞭，忽复落男子之后，而男女之革命，


% 原书第 498 页（PDF 第 511 页）
\BookPage{498}{511}
\noindent 即与种族、政治、经济诸革命并行，成则伸世界惟一之女权，败则同归于尽，永不受制于男。」

在〈公论三则〉中，何震提出三个命题：

一是帝王等同于娼妓。「以一女配无量之男，莫若娼妓；以一男配无量之女，莫若帝王。娼妓无定夫，帝王无定妻；娼妓受污于众男，帝王受污于众女。然而娼妓之贱人知之，帝王之贱人则不知。岂非习惯使之然邪？夫嫖客欲得名妓之欢，犹宠姬欲得帝王之欢也；名妓爱慕嫖客之金，犹帝王爱慕宠姬之色也。两客相遇则争风，犹两姬相值则争宠也。名妓之于嫖客，或疏或亲；帝王之于众妃也，亦或爱或弃。名妓有所欲，嫖客不敢不从；帝王有所命，妃妾亦奉命唯谨。故以男界拟娼妓，娼妓实女界之帝王。以女界拟帝王，帝王实男界之娼妓。况娼妓之翘出者，衣必文采，饰必珠玉，居必洞房，虽帝王之奉，无以加。娼妓舍接客而外，别无所求。帝王舍荒淫而外，亦别无所图。娼妓、帝王，一而已矣。如曰帝王为尊，则娼妓亦为尊；如曰娼妓为贱，则帝王亦可贱。盖女界中之最贱者，以娼妓为最；男界中之最贱者，以帝王为最。欲杀女子，必先杀天下之娼妓；欲诛男子，必先诛亚洲之帝王。况满洲之那拉氏，拥帝王之权，行娼妓之实，帝王、娼妓，毕集于彼一身，此更人人得而诛之者也。」（〈帝王与娼妓〉)

二是政府等同于大盗。「盗一人之所有谓之盗，盗十人、百人、千人、万人之所有，是为大盗。有有形之大盗，有无形之大盗。无形之大盗，政府是也，资本家是也。资本家用攫财之术，以一人之身，而兼有百、千、万人之财。盗百、千、万人之财，而归于一人，下民安得不贫？政府用集权之术，以数人之身，而兼有百、千、万人之权。盗百、千、万人之权而归于数人，下民安得不弱？故下民者，均被盗之人也。政府及资本家，均为盗之人也。下民被盗，穷乏无归，不得不为盗。下民为盗，则罪其身。政府及资本家为盗，所盗百倍于下民。匪唯罪所不及，且为一般社会所尊。世界有此不平之事耶？律以盗律，凡财产增下民一倍者，其罪同于盗一室。若不仅一倍，其罪亦援是为增。凡权力加下民一等者，其罪同于侮一人。若不仅一等，其罪亦援是为增。援此例以讼中国政府，则政府之罪，虽服至重之刑，犹不足以蔽其辜。盖天下无不诛之大盗。政府为大盗之渠首，岂转得赦之而不诛？


% 原书第 499 页（PDF 第 512 页）
\BookPage{499}{512}
\noindent 俗儒不察，犹欲饰儒家尊卑上下之词，以为大盗辩护。吾请诵老子之言，以告之曰：『圣人不死，大盗不止。』则试先废圣人之书，而后治大盗之罪。」（〈政府与大盗〉）

三是权力决定道德。「自古及今，安有所谓道德哉?道德者，权力之变相也。专制之朝，为君者虑臣之背己，又欲臣之为己效死，则以忠君为美德，以叛君为大恶。为夫者，虑妇之背己，又欲妇之为己守节，则以从夫为美德，以背夫为大恶。盖道德者，定于强者之手者也，又强者护身之具也。而道德之效力，则约于权力同。今之恒言，犹以守国家法律为美德。夫今日之法律，何一非强者之法律，迫以实力，使之不敢不从。又诱以虚名，使之不得不从。愚民不寤，又从而和之，曰道德、道德。稍有与之相背者，则斥之为非。夫所谓是非者，强者所定之是非也。强者之对于弱者，凡权力所能制者，制以权力。权力所不能制者，制以道德。权力制人于有形，道德制人于无形。使无量之人，屈服于空理之下，莫敢抗己。强者何其智，弱者何其愚。此真所谓空理杀人矣。试观中国理学之儒所倡诸说，何一而非服从？其尤甚者，则谓君虽不仁，臣不可以不忠。夫虽不贤，妻不可以不贞。天下唯『忠贞』二字，最便于专制之人，非君权、男权盛昌之世，决不至定此讹名。乃腐儒俗子，复从而为之词，合理与势为一谈，即以权力为合于道德，由是权力之外无道德，舍理论势，以势为理，习俗相沿，不以为异。非所谓暗无天日者乎？欲扫荡现世之权力，必先扫荡现世之道德。无论道德、法律，均视为刍狗，则世界之公理，必有复现之一日。处今之世，如有凭现行之道德以决是非者，吾则视为强者之奴。」（〈道德与权力〉）

在〈震致留日女学生书〉中，何震指出:「近岁以来，女界同胞，以婚姻自由、男女平等为恒语，窃以为过矣。昔班昭续《史》，木兰从军，梁夫人援桴击鼓，卫夫人以书法著名，此四贤者，才非不着，名非不闻，身非不正，惜其理不明耳。时婚姻自由、男女平等之说，未克能倡。今女界同胞则不然。悯女学之不兴，则以教育为己任；悲国民之孱弱，则倡尚武之雄风；慨邦国之垂危，则抱治平之雄略，其志可谓壮矣。然予则以为过者，则以日言教育，而无所着之书；日论自强，而不能效身军旅；日议国事，而不闻效力国家。所倡之言论，不过博空谭之名


% 原书第 500 页（PDF 第 513 页）
\BookPage{500}{513}
\noindent 而已；及考其才，亦不过仅留学日本数年而已。此亦何足恃哉？其所恃者，不过以婚姻自专，不告于父母，不宣于宾朋，甚至于数载之间，离合之事，已遭再三矣。此得谓之自由乎？况复致贵重之饰，衣必文采，负粉凝脂，所致之件，已非己力能胜，岂非仰他人之鼻息乎？即不仰他人之鼻息，所欲如此，又安足图大事乎？此非能平等自由之术也。予闻名能传于世界、身朽而人能知者，其人必于天下有实益。今女子受制于男，历有年所。生当斯时，不能大有所作为，以泄女子之耻，以为女界同胞谋幸福，已属可羞之甚矣。顾乃弃立身之术，而取野蛮自由，去实学以空言平等，使闻者生畏，言者掩耳。名誉既失，则男子之压力愈深，非所谓借寇兵而赍盗粮乎？损身以益仇，内自虚而好空言，虽日求自由平等，弗可得也。」

本期杂志还发表了去非子〈破坏社会论〉、不公仇〈李卓吾先生学说〉、志达〈时评四则〉（即〈伟哉女杰〉、〈悲哉男权之专制〉、〈请看俄国虚无党〉、〈女为人妾〉）、公权（即汪公权）〈译稿五则〉（《百北尔总同盟罢工提议案》[译日人《社会主义研究》第五册]、〈巴枯宁学术要旨〉[节译日人《独立评论》第五册]、〈鸣乎劳动者〉[译幸德秋水氏〈平民主义〉、〈击石火节译〉[译幸德秋水<平民主义〉1、〈日本社会歌〉[译意]）、陆守民〈论女子受制之原因〉、大鸿（即周怒涛）〈杂记六则〉（〈女娲为女杰〉、〈异族凌辱女子〉、〈洪秀全男女平等之制〉、〈秦会稽刻石〉、〈郦山女保塞〉）。所列通信所为东京曲町区饭田町六丁目廿一番地何寓

6月25日，《天义报》第二号出版。本期「社说」栏开始刊载何殷震〈女子复仇论〉长文，继续阐发「男子为女子之大敌」的观点，认为「今男子之于女子也，既无一而非虐，则女子之于男子也，亦无一而非仇」。「为男子者，虽受制异族，受制君主，受制资本家，然被治者男子，主治者亦男子也。女子则不然。贵为王后，其身不可谓不尊，而受制于男，自若也。贱为乞丐，其身不可谓不卑，而其受制于男，仍自若也。非惟古代为然，即今代亦然；非惟中国为然，即外邦亦然。「使女子而为人类，又安能日受压抑而不思抵制乎？」

何震认为：「夫满洲之命，固不可不革。然吾则以为汉族之君，其祸更甚于异族之君。何则？汉族之君主，其功愈高，其蹂躏吾女界也愈甚。」「汉族之君主既


% 原书第 501 页（PDF 第 514 页）
\BookPage{501}{514}
\noindent 为女子之敌，故异族为君，其命当革；即汉族为君，其命亦当革。所以革满洲之命者，以其以异族之民，专制吾女界也。且内而政府，外而官吏，均以男子操其权，故满洲之命，应革于吾女子之手。若徒执攘夷之言，附和男子，此何异汉人助满洲人排外耳？且所以攘异族者，为其专制也。专制之政府，固当颠覆；即易专制为立宪，易立宪为共和，然既设政府，则吾人均有颠覆之责。盖政府既设，即有统治机关，而统治机关，必操于男子之手。是与专制何异？即使男女同握政权，然不能人人均握政权也，必有主治、被治之分。以女子受制于男，固属非公。以女子而受制于女，亦属失平。故吾人之目的，必废政府而后已。政府既废，则男与男平权，女与女均势，而男女之间，亦互相平等，岂非世界真公之理乎？」「惟土地、财产，均为公有，使男女无贫富之差，则男子不至饱暖而思淫，女子不至辱身而求食，此亦均平天下之道也。依此法而行，在众生固复其平等之权，在女子亦遂其复仇之愿。盖女子之所争，仅以至公为止境，不必念往昔男子之仇，而使男子受治于女子下也。」

何震详细考察「女子受制之原因」，认为，图腾社会之制，「乃公夫公妻之制度也」，「故血胤相继，咸以女而不以男」，古用女统，「故重男轻女之说未昌」。「及夫两部相争，必此胜而彼败。战胜之族，对于战败之民，施行虐政，惨毒频加。男子则尽遭屠戮，女子则身为累囚。此即劫掠妇女之始也。惟其劫掠妇女，故视女子为至卑，即私女子为己有，备其服役，供其荒淫。是为女子受制于男之始，亦女子属身于男之始也。女子受制于男，历时稍久，遂成习惯，而所掠之女子，又不足以给其欲。由是于同族之妇女，亦出资购买，而所购之女，亦与奴隶齐观，故所生之子，咸易从父姓，而易女统为男统。为女子者，非惟所生子女从父得姓也，即己身亦改从夫姓。」

何震感慨道:「夫子女为父母所共育，岂为父者所得私？且子不母属，为母之人，又何赖于有子乎？今西方各国侈言男女平权，然既嫁之后，均以夫姓自标，即他人亦称以夫姓。夫为夫之人，未尝从妻得姓也，何为妻者竟改从夫姓？岂非男子战胜女子之一大纪念邪？欧美女子号为文明，何改从夫姓，竟习俗相沿，不以为异？吾深惜其愚，吾尤憾其不知耻！


% 原书第 502 页（PDF 第 515 页）
\BookPage{502}{515}
何震认为，「一夫多妻之制，以亚洲为最盛。亚洲各国，又以中国为最甚」。「惟其夫妇不平等，故父母之于子女也，亦不能平等。」「非惟父母重男轻女也，即为子女者，亦重父轻母。」所以，「女子生于中国，少则见轻于父母，长则见轻于夫，老则见轻于子女。」

何震从多个角度分析「男女不平等」现象。「一曰征之于文字。古代之于妇女，不以人类视之，仅以财物视之耳。」「古代之时，女子之名最贱，其义务亦最多。受系累者，女子也；充贡献者，女子也；勤力役而操工作者，亦女子也。男子以女子为生财之具，故得女子愈多者，其家愈富，与南美使役女奴相同。」「以迄于今，女校之中，犹首崇家政。」「夫所谓家政者，非即佐男子治家政乎？佐男子治家政，非古代使役妇女之遗风乎？」「且三代之时，夫人自称曰小童；后世之女，均自称为妾，或以奴家自称，岂非女子甘以奴隶自居耶？」「以妇人与孺子并称，其轻视女子为何如？」「以女子与小人并称，其贱视女子为何如？是虽籍尽女子之口，不能掩压抑妇女之迹也。」

二曰征之于礼制。何震认为，从婚礼制度看，古代有劫掠妇女之风，后世有买卖妇女之俗。今世相传婚礼，均可窥见其遗风。「夫劫妇之俗，由于视女子为俘囚；买女之风，由于视女子为仆隶。今也，仆隶、俘囚之贱，毕集于女子之一身，名曰下嫁，实则以俘囚、仆隶自居。于此而不知自耻，是必女子非人类而后可矣。」从丧礼制度看，「中国之丧服，则男女极不平等」:「重父轻母」，「夫尊妻卑」，「重男轻女」。「推其原因，则夫为妻纲之说有以致之也。」「盖中国之礼制，定于男子之手，故创为此制，以抑女权，安能奉为天经地义乎？」即祭礼、宾礼，亦均失平。「观于礼制之失平，即可知男女之间，自古迄今，均区阶级。蔽以一言，即以女子为为男子附属物耳。」

三曰征之于学术。何震认为，「上古之世，男子以学自私」，「一切之学术，均发明于男子，故三代之书，均含有轻女重男之说」。「秦汉以下之学术，大抵奉儒家为依归。儒家之学术，以重男轻女标其宗。孔丘者，儒家之鼻祖也。而其人即以出妻闻。传子暨孙，均有出妻之行。盖以暴行施之于妻，莫孔门若。」「盖儒家之学术，偏于专制，便于男子之自私，故多妻之说，贞节之风，均儒家有以开其


% 原书第 503 页（PDF 第 516 页）
\BookPage{503}{516}
\noindent 先。汉人祖术儒家，于古籍言及女子者，不惜望文生训，强古经以张己意。」「宋儒继兴，益崇压制，而贱视妇人，屏之人道而外。自是以降，学士大夫莫不执汉、宋二朝之说，奉为玉律金科，习俗相沿，不悟其非。黠者援饰其说以自便，愚者迷信其说而不疑。而吾女子之死于其中者，遂不知凡几。故儒家之学术，均杀人之学术也。」「非扫荡儒书之邪说，则真理无复昌明之期。」

何震认为:「夫男女赋质不同，其为人则一也。」她对班昭《女诫》为代表的传统女学进行了严厉批判，认为《女诫》「首崇卑弱」，「此说一昌，而为女子者，遂以受制于男为定分。名曰礼教，实则羞辱而已；名曰义理，实则无耻而已」。何震直呼其为「贼」，谓「夫班贼身为女子，竟惑于儒家之邪说，自戕同类，以贻女界之羞。作男子之奴隶，为女子之大贼，女界而有此人，盖不啻汉人中之有曾国藩也！」「夫后世女权之不伸，由于为女子者，悉诵班贼之书，以先入之言为主；而班贼之为此言，又由于笃守儒书，以先入之言为主，则班贼之罪，又儒家有以启之也。」

何震认为，古代「惟其贱视女子，故于女子之权利削夺殆尽」，所以女子无兵权、无政权、无学权。「无学权，则女子日愚；无政权，则女子日贱；无兵权，则女子日弱。愚则不能自立，贱则不能自伸，弱则不能自卫。不能自立，则饥饿而死者日多；不能自伸，则屈抑而死者日众；不能自卫，则死于杀戮劫掠者，又不知凡几。岂非古今之男子，无一非迫女子于死境乎？虽然，此犹男子用间接之法以死吾女子者也。其以直接之法死我女子者。」「一曰死于幽闭，二曰死于羞辱，三曰死于刑戮，四曰死于劫掠，五曰死于摈弃，六曰死于贞节。」何震对这六个方面进行了深入的剖析。

本期还发表了〈法国女杰露伊斯米索尔像〉、亚公〈唐铸万先生学说〉、大鸿〈时评四则〉（〈清政府亦欲抑女权耶〉、〈异哉中国妇人会〉、〈雇婢评〉、〈社会革命水风潮〉）、大鸿〈杂记六则〉（〈女娲为女杰〉、〈赵威后语〉、〈异族凌辱女子〉、〈洪秀全男女平等之制〉、〈秦会稽刻石〉、〈郦山女保塞〉）、公权〈译稿三则〉（〈法国女杰露依斯传〉[意译日人志津野又郎之作]、〈俄国女杰骂阁臣语〉[节译日人〈新时代报〉]、〈日本劳动歌译意〉）、恢权〈平


% 原书第 504 页（PDF 第 517 页）
\BookPage{504}{517}
\noindent 权论〉、留学界一分子〈东京来函〉、白圭女史〈女地狱歌〉、叔时〈杂记三则》（〈中国古代均贫富之制〉、〈富人之病民〉、〈颜元斥土地不均之害〉）。另有〈本社捐金芳名〉共列四人:黄中黄君五元，顾宝瑚君五元，沙浮君三十元，周大鸿女士五元。

本月，日本幸德秋水读《天义》第一期后，既赞赏何震等的革命精神，同时又对其中的观点提出质疑，于是致书何震，邀其面商，信函称：「《天义报》首册读毕，拍案呼快。就中〈女子宣布书〉，议论雄大，如名将行兵，旗鼓堂堂不可当；若〈帝王与娼妓〉，骂得痛绝，如利刃刺骨，何等刺心文字也！敬服之至。男女同权者，人生之精理，而方今之急务也。贵娘今开阐此真理，以冀实行此急务，以着女子先鞭，感激无已；即为社会之女子思之，固不容不谢此劳也。但仆意犹有憾者，则以仆所持者，有一二之异见，可容仆之直言乎？贵娘所言，以为初婚之女，必嫁初婚之男；再婚之妇，必嫁再婚之夫。是诚仆之所不解也。夫夫妇关系之第一要件，在于男女相恋相爱之情。纵令初婚之夫妇，心中无相恋相爱之情，则固有妨于夫妇之道。又令再婚之男与初婚之女，真克爱恋和谐，何害其为夫妇乎？而贵娘必欲使初婚之男女、再婚之男女互相配合，能无仍为古来『贞女不见二夫』之陋道德所染乎？仆窃疑之。既以爱情为男女交际之要件，即可不必问其为法律上所许否，即夫妇相栖而成家，如现时所谓结婚制度者，将来自由之社会，亦不必以此为必要之事。盖仆最恶淫佚，为于伪道德之束缚人生爱情者，更觉其有害。则男女问题之解决，岂可不深思乎？仆今卧病褥，艰于执笔，尚幸贵娘枉驾茅屋。又，友人堺利彦，平生颇讨究女权问题。贵娘若来，仆亦欲延之共论也。」

何震等即往拜访，晤幸德秋水、堺利彦等，交换看法。并为《天义》第三卷刊载的〈幸德秋水来函〉加案语，说：「震得书后，即往访幸德君，时堺君亦在座。堺君所持之见，略与幸德君函中所言相同。盖幸德君及堺君之意，在于实行人类完全之自由；而震意则在实行人类完全之平等，立说之点，稍有不同。即幸德君及堺君，亦以为持论不必强同也。」

7月10日，《天义》第三卷出版，「图画」栏刊有〈露国革命志士之气概〉、<名画之二〉（曼殊作〈邓太妙秋思图〉），「社说」栏续刊震述〈女子复仇


% 原书第 505 页（PDF 第 518 页）
\BookPage{505}{518}
\noindent 论〉，刊载申叔人类均力论〉（内页题作「人类均力说」）；「学理」栏刊去非〈戴东原先生学说〉；「时评」栏刊载志达的〈政府者万恶之源也〉和〈排满与保满〉；「译丛」栏刊载公权〈苦鲁巴特金之特色〉（内页题下注:「此摘译日本久津见蕨村《无政府主义》」）、〈妇人与政治〉（题下注:「译日人幸德秋水平民主义〉）；「来稿」栏刊载〈幸德秋水来函〉；「附录」栏刊恢权〈女子军歌〉、震述〈陈君不浮追悼会演说稿〉；「书报绍介」栏介绍了日文书刊神崎顺一〈世界妇人报〉、堺利彦〈妇人问题〉和〈家庭杂志〉、〈熊本新闻〉、西川光次郎着富之压制〉、堺利彦译〈理想乡〉，英文书籍金子喜一夫人编辑〈社会主义妇人杂志〉。所列通信所为东京牛込区新小川町二丁目八番地何寓。

何震在〈陈君不浮追悼会演说稿〉中指出:「今天所开的会，是因为追悼陈君不浮。陈君的死，是因为这一次回国，听见人讲，说现在福建的地方，已经归日本势力范围，所以愤不顾生，竟投海而死。我对于这件的事情，实在是恨极了强权的说话。甚么叫做强权？就是强盗的别名。同是个人类，凡权力强的，人人可以得利；凡权力弱的，人人都要受祸。他外面说的话，还要说强者对于弱者，应该蹂躏。这真是不道德的极点了。所以，实行强权的人，就是杀人抢钱的人。现在的世界，都讲帝国主义，都讲强权。可以直接杀人，如法国待安南，英国待印度，都是用这种法子的；又可以间接杀人，如这次陈君蹈海，因为日本想侵略福建，这就是受强权影响而死的。足见强权这句话，无一不足以杀人。况且，强权的学术，既然盛行，所以同国的人，也都是扶强锄弱。这一种思想，正与中国人的势利心相合，所以现在的中国，都是只论势力，不论是非，仿佛『势力』两个字，所含的意思，同『理』字一样。试想:天下的人，都是人类，为甚么强者独生、弱者独死呢？为甚么强者独利、弱者独害呢？照中国人心这样行，必至于一个人有利，社会都受他大害。照欧洲、日本人心这样行，必至于一个国有利，世界都受他大害。所以，我对于陈君蹈海，到格外恨强权。但现在的人，所发的理想，不过到抵制强权为止，不晓得抵制强权，还是以暴易暴。想抵制强权，不如打消强权。想打消强权，都要扶弱锄强。对于力强的人，无一不用其抵制；对于力弱的人，无一不用其保护，教强弱复归于平等。如若不归於乎等，惟有用陈君遗书所讲的话，人人不怕死，人人不顾害，把一切不公不平的社会，都由暴力破坏。大约现在的社会，都是阶级社


% 原书第 506 页（PDF 第 519 页）
\BookPage{506}{519}
\noindent 会。阶级社会，就是强者的社会。现在的世界，都是功利的世界。功利世界，也是强者的世界。不破坏社会，决不能打消强权，决不能教人类平等。所以，诸君对于陈君蹈海，当晓得陈君蹈海的原因。陈君蹈海的原因，是因为不忍见强权盛行。我望诸君因陈君这一死，个个人明白强权的流毒，这就不负陈君一死的苦心了。」

7月6日徐锡麟在安庆枪杀安徽巡抚恩铭，率领巡警学堂学生起义，兵败被杀；14日，秋瑾因胡道南告密，次日在绍兴就义。

25日，《天义》第四卷出版，其「图画」栏刊载〈烈士徐君锡麟陈君伯平遗像〉、海哥美尔作〈同盟罢工〉（名画之三）、曼殊作〈江干萧寺〉（名画之四），「社说」栏除续刊何殷震〈女子复仇论〉外，还刊发了申叔〈无政府主义平等观〉，「学理」栏刊申叔〈西汉社会主义学发达考〉，「时评」栏刊发志达〈时评四则〉(〈不平哉万国平和会〉、〈暗杀之影响〉、〈鸣乎三纲之国〉、《专制之政将复活矣》），「译丛」栏刊登了公权〈西欧无政府党宗旨汇录〉和〈婆利箫斯楷传〉，「来稿」栏刊载汉一毁家论〉、独应（即周树人、周作人兄弟）〈妇女选举权问题〉和绝诗三首（别女界也）〉。据本卷〈广告〉称:「独应君鉴：前承惠寄大作，不胜铭感。惟未详尊寓所在，乞急示知。因有要事奉商故耳。」则此时何震等与周氏兄弟尚不相识。本卷所载〈本社捐金芳名〉计列四人，分别是：曼殊君五十元，汪公权君四十元，周大鸿女士四十元，陆恢权女士二十元。

8月10日，《天义》第五卷出版，其「图画」栏刊载〈中国无政府主义发明家老子像〉、曼殊〈清秋弦月〉（名画之五）；「社说」栏除续刊何殷震〈子女复仇论〉、申叔〈无政府主义之平等观〉外，新刊畏公〈论女子劳动问题〉；「学理」栏续刊申叔〈西汉社会主义学发达考〉，「时评」栏刊志达〈时评二则〉（〈男盗女娼之上海〉、〈政府奖励实业〉），「译丛」栏刊去非〈杜尔斯德答日本报知新闻社书〉；「来稿」栏刊悲生〈秋瑾传〉、太炎〈秋女士诗集序〉、子谷〈秋瑾遗诗序〉，「附录」栏刊何震编〈曼殊画谱序〉（河合氏、章炳麟、曼殊、何震），〈新刊绍介〉（《新世纪》杂志）。所列通信所为东京小石川区久坚町二十七番地何寓。

何震〈曼殊画谱后序〉云:「古人谓境能役心，而不知心能造境，境由心而


% 原书第 507 页（PDF 第 520 页）
\BookPage{507}{520}
\noindent 生。心之用无穷，则所造之境亦无极。如绘画一端，古代皆以写象为工，后世始有白描山水，以传神擅长。其所以易写象为传神者，则写象属于唯物，而传神近于唯心。画而出于白描，此即境由心造之证也。吾师于唯心之旨，既窥其深，析理之余，兼精绘事；而所作之画，则大抵以心造境，于神韵为尤长。举是而推，则三界万物，均由意识（构）[构]造而成。彼画中之景，特意识所构之境，见之缣素者耳。此画学与惟心论相表里者也。因汇为《画谱》，先将第一集开印，余俟续出，并乞吾太师母、吾师及太炎先生序而行之，以问十方高士。丁未初秋，仪征何震手书」

8月31日下午一时，社会主义讲习会在东京牛込赤城元町清风亭举行第一次集会，何震参加会议并发表演说，称：「吾于一切学术，均甚怀疑，惟迷信无政府主义，故创办《天义报》，一面言男女平等，一面言无政府。盖无政府之目的，在于人类平等及人无特权。若男女平等，亦系人类平等之一端；女子争平等，亦系抵抗特权之一端，并非二主义相背也。特无政府主义，不仅恃空言也，尤重实行。现世界无政府党，以俄国为最盛。俄国无政府党，其进步分三时期：一为言论时代，二为运动时代，三为暗杀时代。今中国欲实行无政府，于以上三事，均宜同时并做。即使同志无多，亦可依个人意志而行，以实行暗杀。盖今日欲行无政府革命，必以暗杀为首务也。」（《天义》第六卷，公权〈社会主义讲习会第一次开会纪事〉）这是何震在讲习会历次活动中唯一一次参会和公开演讲。刘师培、张继、幸德秋水也在会上发表演说。

9月1日，《天义》第六卷出版，其「图画」栏进载〈徐烈士（徐锡麟）遗墨〉、法民罢工传单写真〉，「社说」栏续刊畏公〈女子劳动问题〉，刊载震、申叔合撰〈论种族革命与无政府革命之得失〉，「学理」栏刊载申叔〈欧洲社会主义与无政府主义异同考（驳鹤卷町某君来函）〉，「记事」栏刊载公权〈万国社会党大会记〉、〈记日本社会党欢迎英开革哈迭氏记事〉、〈社会主义讲习会第一次开会记事〉、〈日人社会主义金曜讲演会广告〉，「来稿」栏刊载〈绍兴某君来函论秋瑾事〉，〈新书介绍〉（〈母亲之土地〉、〈结姻离婚〉，「附录」栏刊载何震录〈梵文典序〉（章太炎序、刘师培序、熙州仲子题诗、何震题偈）。柳亚子


% 原书第 508 页（PDF 第 521 页）
\BookPage{508}{521}
\noindent 〈曼殊新传〉云：「丁未，从章炳麟、刘师培游，着《梵文典》八卷，自为序。师培为《天义报》，倡无政府主义，邀玄瑛同居，刊其画于报端，师培妇何震则从玄瑛习绘事，号称女弟子。震为玄瑛辑画谱，玄瑛自有序，河合氏暨炳麟为序，震为后序，将付梨枣。又思刊布《梵文典》，印度波逻罕学士暨炳麟、师培为序，震为题偈，陈独秀为题诗，顾二书咸未及刊成，仅于《天义报》载其序跋而已。」

何震、刘申叔〈欧洲社会主义与无政府主义异同考〉认为，「社会主义多与无政府主义相表里。考无政府之学术，本分二派。一以个人为主」，「以个人之发达进步，期无政府主义之实现」「一以社会为主」，「以社会经济之改革，期无政府主义之出现」。「此二派者，一主吾人内部之改革，一主吾人外部之改革；一主心意、性格之发达，一主境遇、事情之变迁；一主理论，一主实行。盖一为哲学的无政府主义，一为合群的无政府主义也。哲学的无政府主义，以个人主权为重，诚与社会主义不同。若合群的无政府主义，其共产一派，则与社会主义不殊；其不同之点，则一欲利用国家，一欲废绝国家耳。夫社会主义之目的，在于冀一切财产之平均，以易私有为公有；若共产的无政府主义，其目的正与此同。」「社会主义以平等为归，既有政府，则必上有统治之人，下有分配之机关，均背于平等之旨。故由社会主义扩张之，必达无政府主义之一境；而共产的无政府主义，实由社会主义而生，乃社会主义之极端，不得别之于社会主义之外也。况社会主义欲借国家之力，以均平财产者也，然亦易为国家所利用。昔法国露易佛兰，于国立工厂斥为伪社会主义。今倡土地财产国有之说者，说非不善，然不善用之，必致一切利权，均为国家所垄断，以蹈法国国立工厂之覆辙。若行无政府主义，则此弊无自而生。」「若徒执社会主义而排斥无政府主义，此则偏于一隅之见耳，非吾等之所欲出也。」「社会主义之实行，以劳动集合为嚆矢。欲行无政府主义，亦以劳动集合为权舆。盖欲行无政府革命，而不能使多数平民均明无政府主义，则其制不能实行。故吾人欲行无政府，亦从劳动集合入门。其宗旨虽与社会党不同，至其行事，则固未尝迥殊也。」

<绍兴某君来函论秋瑾事〉案语称:「秋瑾此次之惨死，为今古莫大之冤狱。吾非谓秋瑾非革命党员，特以绍守之杀瑾，既未得其确实之供，复伪造供词，以淆


% 原书第 509 页（PDF 第 522 页）
\BookPage{509}{522}
\noindent 乱事实，此其所以为冤也。吾闻此次绍兴狱事，系因袁翼告密。袁翼者，与锡龄同为某暗杀团党员。及蒋观云创政闻社，鼓吹立宪邪说，又引翼为己党。观云素与瑾相识，瑾固多大言，尝语人曰:『蒋观云者，吾司为东京革命机关。』蒋恐祸及己，恨瑾刺骨，故与袁翼通谋，促之使告密。及东京绍兴人开同乡会，观云犹语人曰:『为保卫地方上治安计，不得不杀瑾。』其设心之毒如此。若夫所以成此狱者，则为浙江抚幕张某及汤寿潜。汤寿潜者，外托伪道德之名，隐为谋利之计，为浙江之大贼。此次绍兴之狱，系因寿潜运动浙抚、绍守。浙抚、绍守信其言，故成此狱。近寿潜致书东京某同乡，盛夸绍守办事之善，阅者皆为目裂。故吾于秋瑾之冤，不怨浙省官吏，惟怨浙省士大夫。」

何震〈梵文典题偈〉云:「佛说一切法，聚十方僧众，运此广长知，说无尽妙言。一佛一舌舌一音，以为一切语言海。试从声闻参正觉，普放光明证自在。弟子何震拜题。」

本卷所载〈本社捐助及赞成芳名〉计十三人，分别是：唐群英女士贰元，功群女士伍元，方君英女士贰元，杨蕙畹女士壹元，何香凝女士六元，丘新荣女士伍元，廖平女士贰元，杨光君贰元，邱心荣君壹元，方萌女士贰元，王寄庐君拾元，郭君女士伍元，胡汉子女士三元。其〈本社广告〉称:「本报从本册起另增记事一门。又本册因所记之事浩繁，暂缺时评、译丛二门。」又称：「本报自第一册起本定于每月初十、廿五二日出版，现因本册出版逾期六日，本册以后定于每月初一、十五二日出版，以昭画一。」

秋瑾逝世后，何震潜心收集其诗词，编成《秋瑾诗词录》二卷，计收录诗八十七首，词三十八関，各为一卷。该由王芷馥助资刊印于日本东京，平装四十九页，为秋瑾就义后出版的第一个作品集。

9月2日，何震作〈秋瑾诗词录后序〉:「秋瑾罹祸之岁，七月初旬，得其诗词若于首，各为一卷，乞太炎先生及吾师曼殊为序，并由吾友王芷馥女士助资排比。阅二旬而成，乃书其后曰:人治者，摧折天才之具也。天才贵肆，人治贵拘。惟尽决人治之藩，斯天才之发育依自然而舒，充然有以自遂；否则，灵智虽具，遏于人治，亦斲其萌蘖，以归于梏亡。如中国妇学，兴于三代之前，而自古迄今，才智之


% 原书第 510 页（PDF 第 523 页）
\BookPage{510}{523}
\noindent 女不数出，其故何哉？盖妇教设官，首崇四德，而公宫女吏之所戒，其条目节次，均立文以为坊。由是足不逾阈，跬伏闺闼，安于所习，蔽所不见，以锢其思索，障其灵源；学业无所证，志气亦鲜所发舒，胶固沉溺，有若天囚。虽以诗言志，而诗之本真不见。旷观古今女子，其赋诗倚声者以十百计，然标托引喻，伊郁善感；诗词弥工，而托体弥卑。会稽章氏谓：『古之女子，因习礼以达诗。』不知古之所谓礼者，揭人治以为范，使絷身礼教之中，则高旷之志，灵敏之质，均将汩没于无形，以反于顽冥。处秉礼之世，安有所谓性灵之诗哉？是则以人治遏天才，于女子为特甚。秋瑾者，女子之富于天才者也。反古易常，不为纲维所域；又执持光复之谊，谙悉清廷政教之非。尝再莅日本，往返楚越燕云间，以与志士相结纳。虽言行不自检，然尚气节，重然诺，大昌侠烈之风，均与古代妇学异轨。今读其诗词，多慷慨之音。凡欢愉忧愤之情，身世家国之感，一寄之吟咏。思有所寄，援笔直摅，而生平志节，又隐约于意言之表，殆古之所谓性灵之诗欤！虽然，瑾之克舒其才，瑾之不囿于人治也。傥人治既废，则女子之具天才者，不为人治所束，势必益臻于灵智，其心知之瀹，必非今日所能跻。然则今之兴女教者，其发抒其才，以遂其自由之性乎，抑将囿于浅俗以遏其灵敏之机乎？观于秋瑾之诗词，可以判其得失矣。丁未年七月二十五日，何震识。」

9月7日上午，何震组织召开追悼会，悼念徐锡麟、陈伯平、马宗汉和秋瑾。《天义》第七卷〈徐陈马三烈士及秋女士追悼会记〉称：「本报编辑人何女士震，于西历九月七日，开徐（锡麟）、陈（伯平）、马（宗汉）三烈士及秋（瑾）女士追悼会于神田锦辉馆，会者千余人，遂于午前九时开会。先由何女士报告开会，次由章君炳麟报告徐、陈、马、秋之历史，次由刘君光汉演说暗杀之效果。相继演说者络驿不绝。次由张君继读祭文（祭文系章君太炎所作，见《民报》中）。次行三鞠躬礼。至午后三时散会。」

9月15日，《天义》第七卷出版，其「图画」栏刊载〈维拉百态哥奈尔〉、罗马教之看护妇（博爱主义）〉照片，「社说」栏除续载申叔〈无政府主义之平等观〉和震、申叔合撰〈论种族革命与无政府革命之得失〉二文外，刊载了震述<女子解放问题〉，「学理」栏刊载民鸣〈社会契约说评论〉，「时评」栏刊载志


% 原书第 511 页（PDF 第 524 页）
\BookPage{511}{524}
\noindent 达〈时评三则〉（〈结婚奇谈〉、〈世界主义之公敌〉、〈亚洲之暗杀〉），「记事」栏刊载公权〈记事七则〉（〈万国无政府党大会记略〉、〈伊大利社会党分裂记〉、〈徐陈马三烈士及秋女士追悼会记〉、〈社会主义讲习会预告〉、〈日本社会主义金曜讲演会记〉、〈清外部欲禁新世纪不果〉），「译丛」栏刊载独应〈妇人选举问题〉，「杂记」栏刊载独应〈读书杂拾〉，「杂录」栏刊载何震〈秋瑾诗词后序〉。所附〈本社捐助及赞成芳名〉计列五人：陈登山君五元、张旭庐君五元、李萼女士伍元、周锡之君五元、李季恭君五元。

在〈女子解放问题〉中，何震认为:「数千年之世界，人治之世界也，阶级制度之世界也，故世界为男子专有之世界。今欲矫其弊，必尽废人治，实行人类平等，使世界为男女共有之世界。欲达此目的，必自女子解放始。」「中国数千年之制度，以女子为奴隶者也，强女子以服从者也。」「女子毕身之责任，不外育子及治家二端。」「由是，奴隶、惰民之恶，悉集于女子之一身。然为男子者，亦安之如素。」「至于近日，不独女子失其自由也，即男子亦以家室之累，而失其自由。凡奉母、蓄妻、嫁女之费，毕集于男子之身。」「然男子虽躬罹其苦，仍以囿于礼法之故，以解放妇人为大戒。」然而「幽闭女子之制，决不足以禁女子之不淫。夫幽闭女子之目的，既在于防止淫泆，而其结果，势必无解放女子之名，而女子之心，转人人抱一淫泆之念。名曰禁淫，实则诲淫而已。盖既以解放妇人为大戒，又虑解放以后，妇人即从事宣淫，防之愈严，则妇女逾防之念日切，稍有解放之隙，则淫迭之念生。是犹禁人以盗物，为盗者知物之可贵，而盗物之心益切也。故女子之犯奸淫，由于幽闭，而非由于解放，安得谓解放女子即系导女子以淫泆哉？乃中国人民不察其由，愈以解放妇女为戒，此女德所由日堕，而女性所由不发展也。」

「中国之婚姻，礼法之婚姻也。若欧美诸国，则昔日之婚姻，为宗教婚姻之制；近世之婚姻，为法律婚姻之制。其制之胜于中国者，一则结婚离婚，均可自由，兼可再嫁；二则行一夫一妻之制；三则男女同受教育，男女同入交际场。就表面观之，不可谓非解放女子也。然吾谓此等之制，仍属肉体上之解放，非复精神上之解放。何则？解放者，不受缚束之谓也。今观欧美婚姻之制，一缚于权利，再缚于道德，三缚于法律。」


% 原书第 512 页（PDF 第 525 页）
\BookPage{512}{525}
「至于男女平等，则亦弗然。夫男女虽同受教育，然处人治盛昌之世，政治法律，女子攻者甚鲜，而陆军、员警之学，不复令女子与闻。男女虽同入交际场，然处政府擅权之世，官吏之职，不加于女子之身。则所谓男女平等者，有其名而无其实者也。夫解放女子，必使为女子者，共用平等自由之乐。若如今日之欧美之制，势必女子有自由之名而无自由之实，有平等之名而无平等之实。其所谓自由者，非纯正自由也，伪自由耳！其所谓平等者，亦非纯正平等也，伪平等耳！无自由之实，故女性未克发展；无平等之实，故人权未克均平。亚洲妇女，震于欧美之文明，以为欧美女子实行解放，实享平等自由之乐，一若克步欧美女子之后尘，为愿已足。呜乎！处今日女子革命之时代，吾决不望女子仅获伪自由、伪平等也，吾尤望女子取获真自由、真平等也！」

\begin{BookQuote}
「近岁以来，中国之社会，亦渐谋女子之解放。然女子之解放，有真出于主动
\end{BookQuote}

者，亦有出于被动者。」而男子提倡女子解放，一是「因求名而解放女子者也」，二是「因求利而解放女子者也」三是「因求自逸而解放女子者也」。所以「今日之解放妇人，出于男子之自私自利。名曰助女子以独立，导女子以文明，然予女子以解放之空名，而使女子日趋于劳苦。」「男子之解放妇人，亦利用解放，非真欲授权于女。」

何震强调：「吾非谓世界一切之职务，当专属于男，不当为女子所分担，亦非谓女权不当扩张。特以女子之职务，当由女子之自担，不当出于男子之强迫；女权之伸，当由女子抗争，不当出于男子之付予。若所担责务由男子强迫，是为失己身之自由；所得之权由男子付予，是为仰男子鼻息。名为解放，实则解放之权属于他人，不过为男子所利用，而终为其附属物而已。故吾谓女子欲获解放之幸福，必由女子之自求，决不以解放望之男子。」所以，「女子欲真求解放之幸福，正宜发展其女性，以握改造社会之权。若徒知寄情淫欲，则救世之心，或为纵淫之心所夺，所抱之志必不克成。」

针对当时蓬勃发展的「争取女性选举权」运动，何震指出，「夫吾等所谓男女平等者，非惟使男子不压抑女子已也；欲使男子不受制于男，女子不受制于女，斯为人人平等」。因此，「少数女子握权，决不足以救多数女子」。「盖人治一日不


% 原书第 513 页（PDF 第 526 页）
\BookPage{513}{526}
\noindent 废，权力所在之地，即压制所生之地也。」「故为多数女子计，苟非行根本改革，使人人平等，宁舍选举权而勿争，慎勿助少数女子，俾之争获参政权。盖昔日压制多数妇女者，一为政府，一为男子；今则政府及男子而外，另受制于上级之妇人，则是于己身之上，别增一重之压抑也。即使压抑不增，亦仅供少数妇人所利用，夫何幸福之有哉！夫何解放之有哉！」

何震认为:「妇人解放问题，当使为妇人者，人人同享解放之乐。今之持解放说者，一曰女子职业之独立，二曰男女参政权之平等。不知所谓职业独立者属于个人，抑属于全体？如曰属于个人，则仅己身不受制，非多数妇人均可免厄也；如曰属于全体，则以今日经济界之组织，少数富民垄断生产之机关，平民失业，其数益增，而谓妇女职业均能独立，则所谓职业独立者，即以职业供役于人之异名耳。自由、解放，岂可得哉！故谓职业独立则女子可以解放，不若谓实行共产，妇女斯可解放也。至于与男子均权，无论男子握权历时已久，男女参政之柄，非仓卒所能均。即使能均，决不能人人而参政。以少数参政之女子处于主治之位，使多数无权之女子受其统治，不独男女不平等，即女界之中，亦生不平等之阶级。彼多数妇女，不甘受制男子者，岂转甘受制女子乎？故今日之女子，与其对男子争权，不若尽复人治，迫男子尽去其特权，退与女平，使世界无受制之女，亦无受制之男，夫是之为解放女子！夫是之为根本改革！奚必恃国会政策，以争获选举权为止境哉？」

本月中旬，刘师培母亲李汝东渡日本，与何震、刘师培同住。

10月30日，《天义》第八至十卷合册出版。其「图画」刊载〈克鲁巴特金像〉、〈露依斯米索尔画像〉、〈俄冬宫图〉、〈梁萃轩像〉、〈从军苦图〉，「论说」续刊震〈女子解放问题〉、〈女子复仇论〉，刊载申叔〈论新政为病民之根〉、申叔〈中国民生问题论一（野业与市业）〉、申叔〈非六子论〉、震〈论女子当知共产主义〉，「学理」栏刊载申叔〈鲍生学术发微〉、自由〈斯撒奈尔无政府主义述略〉，「时评」栏刊载扶弱〈祝日本社会党之分裂〉、〈惨哉工女〉、〈哀我农人〉，「记事」栏刊载公权《芜湖万顷湖农民抗租记》、〈中国各省罢市毁学汇志〉、公权〈万国社会党大会再记〉、〈无政府党大会决议案记〉、〈记女


% 原书第 514 页（PDF 第 527 页）
\BookPage{514}{527}
\noindent 界与万国社会党大会之关系〉、〈女界近事记〉、〈记东京金曜讲演会〉、〈社会主义讲习会记事并报告〉，「来稿」栏刊载凤池（正文作「凤词」）〈节妇辨〉、少年〈天义之证定〉、〈汉川陈志淑女士来函〉、〈东京五六居士来函〉、〈神田某君来函〉、〈日本小幡宗海来函〉、〈日本长肇君来函〉、〈王朗鸿菱花将军行〉，「译丛」栏刊载秋心室主人〈秋心室丛谈〉（〈伦敦之贫民窟〉、〈人种的自杀〉、〈功利主义〉、〈美洲社会党之二派〉）、衡民俄国女杰遗事汇译〉、「杂记」栏刊载独应〈读书杂拾〉、志达〈女界吁天录〉、申叔〈穷民俗谚录〉、志达〈石破天惊谈〉，「专件」栏刊载〈农民疾苦调查会章程〉，「附录」栏刊载〈南非洲杜省虐待华侨之惨状〉、申叔〈总同盟罢工论序〉、刘申叔〈穷民俗谚录征材启〉、申叔〈活地狱序〉、〈戒学政法歌〉、〈平民歌〉、〈新刊绍介》（〈总同盟罢工〉）。所附〈本社捐助及赞成芳名〉列五人：郑君、黄君合捐二十元，刘三君十元，剑公君十元，汉一君五元。

何震在〈论女子当知共产主义〉中以白话的形式，向穷苦姐妹们作宣传，说：「天下是甚么东西顶要紧？就是吃饭要紧。你们做女人的，甚么要受人的虐待？就是要靠人吃饭了。试看现在顶可怜的女人，共总有三种。一种是做老妈、丫头的。主人要打就打，要骂就骂，一点儿不敢抵抗。还要一天到晚的替他做事，起五更，睡半夜。这是甚么道理？就是因为主人有钱，我想靠他吃饭了。一种是做女工的。上海各处地方，丝厂、纱厂、织布厂、洗表作里面，雇的女工，不晓得多少，也是一天做到晚的，也是一刻不能随便的。弄得眼睛看不见，弯腰驼背。这是甚么道理？就是因为厂里有钱，我想靠他吃饭了。一种当婊子的。天天要被龟头打，无论甚么人，要嫖就嫖，要叫局就叫局，被人看得很轻。还有上海的野鸡，三更半夜的时候，就是大风大雪，还要站在街上等客人。这是甚么道理？就是因为家里没有钱，把自己拿出来卖，靠这种事情吃饭了。这个三种人以外，还有做人妾的，无论受正室这样欺凌，也只好忍气，也是因为靠男子吃饭了。又有做寡妇的，有钱的人家，殉节的很少。到了一种没有钱的人家，又没有儿女，又没能再嫁，死的也就很多。这也是因为没有饭吃。就是有饭吃，过的日子也很苦，所以就自己寻死了。就是种田、养蚕的女人，他也是非常的苦。他甚么一定要做呢？也是因为混饭吃。又如，既嫁男人以后，也有为男子打骂的，也有为男子不理的。他甚么不敢大闹呢？


% 原书第 515 页（PDF 第 528 页）
\BookPage{515}{528}
\noindent 并不是看男人的面子，实在是看饭碗的面子。所以做女子的，因为想这碗饭吃，不晓得受了多少苦，也不晓得吃了多少亏。你们做女子的，不要恨男子，只要恨没有饭吃。为甚么没有饭吃？就是因为没有钱，不能买饭。为甚么没有钱？就是因为有钱的人，把财产掠夺了去，所以弄得多数的人，穷的没有饭吃。你看各衙门、各公馆里面，做太太、做小姐的，何等阔气？他也不愁没饭吃，为甚么你们天天都愁饿死？穷的也是人，富的也是人，你们自己想想，也应该生出一点不平的心出来了。

现在有一种人，说做女人的，只要有一件行业可做，就不怕没有饭吃。譬如中等的人家，把女儿送进学堂，学一点儿普通学，或是学一点儿手工，就是嫁人以后，也可以出去做教习，不至靠男人过活。又如那种很穷的人家，无论女儿、媳妇，把他送到工厂里做工，好在工厂一天多一天，不怕没有安身的地方，也不至再当老妈、丫头、婊子了。这一句话，也有点儿道理。但就我看起来，学堂是人家出钱办的，到里面做教习，就是靠开学堂的人吃饭；工厂也是人家出钱办的，到里面做工，就是靠开工厂的人吃饭。既然靠他吃饭，就一点自由都没有了，共从前靠男人吃饭受他的压制，也差不多。哪里可以叫独立呢？况且靠学堂、工厂吃饭，如若学堂、工厂关闭，或是他嫌人多、被辞了出去，或是才具没有人要，可不是还是没有饭吃呢？所以靠人吃饭一件事，说到末了，都是非常的危险，决不是个顶高、顶上的好法子，就是俗语说的『脚面上支锅』了。

我现今有一个好法子，叫你们不要靠人，自然就有饭吃。这是甚么法子呢？就是实行共产。你想：天地间的东西，不是天生出来的，就是各人做出来的。为甚么有钱就可以买，没钱就不能买？就是因为世界上用钱的缘故，就是因为各人把钱买的东西当作私有的缘故。如若做女人的，个个人晓得钱这样东西，是个顶坏不过的，大家齐心起来，同男人家合在一起，把在上的人同有钱的人，一律废尽，随后把钱这样东西，也废了不用，无论甚么东西，都不准各人私有，凡吃的、穿的、用的，都摆在一个地方，无论男人、女人，只要做一点工，要哪样就有哪样，要多少就有多少，同海里挑水一般，这就叫做共产制度。到那个时候，不独吃饭不要靠人，还天天都有好饭吃，还可以有好的穿、好的用、好的顽，你们想想：这个日子，好过不好过呢？况且，我这句话，并不是骗你们，只要大家齐心，照着这个法


% 原书第 516 页（PDF 第 529 页）
\BookPage{516}{529}
\noindent 子做，自然有这个好日子教你过。你们不必疑惑。俗语说道:『好日子在后头。』就是今天所讲的这句话了。」

本卷还附录了〈社会主义讲习会支金报告〉，至本日，社会主义讲习会「入款项」共十八元，其中：怪汉子五元（犹有十元未缴），刘君五元，许君一元，刘光汉君三元，何震女士四元。「出款项」共十七元二角五分，其中:首次租会场费三元，二次租会场费一元八角，三次租会场费一元八角，四次租会场费一元五角，首次至四次「邮片费」五元一角五分，幸德君演说稿印费四元。存金七角五分。这是社会主义讲习会的第一次财务收支公开报告。所列通信所为东京曲町区饭田町六廿一何寓。

11月上旬，由于旅日开销及办报所费颇大，经济拮据，何震由日本回到上海，逗留数日。受章太炎之托，何震回国后，利用其兄何家辂（誉生）与张之洞女婿卞昌的朋友关系，为章太炎赴印度研究佛经，向张之洞运动路费。何震到上海后，章太炎曾给何震寄去一封信，称:「连接数信，因懒未覆。近想安抵沪上。天寒气燥，宜多吃芦服、橄榄以自卫。所托诸事，务望尽力。曼殊师前有信来，知其有意东渡，欣慰之至。附一函，可转见则转交之。上海传言：兄与申弟皆已入日本籍。此等妄语，不知何人编造，甚可怪也。兄麟顿首。」章刘反目后，何震对外公布此间通信，并加注语，云：「此信无甚关系，惟说『所托诸事，务望尽力』二言，则凡运动张之洞诸事，皆包括其中矣。」

案，卞昌（1873-1946），原名纶昌，字经甫，号薇阁，晚号猥盦，江苏仪征人，居扬州。祖父卞士云是嘉庆三年的进士，历官监察御史、湖北按察使、浙江布政使署巡，着有《退思斋诗集》。父亲卞宝第是同治年间的举人，历任闽抚、湘抚、闽浙总督兼福建船政。卞昌36岁中优行廪贡生，初娶广东廉州府知府张丙炎之四女。光绪二十五年（1899）续娶张之洞长女张仁准为妻。光绪二十七年（1901），出使日本大臣蔡钧奏调他充任随员，派充日本横滨领事、长崎正领事官。卞绰昌在外为使馆官员，在内承办新式企业，得当时风气之先。辛亥革命后归隐不出，优游林下。善书法。卒后葬于仪征刘家集塘坝祖莹。

11月15日，何震由上海返回扬州。25日，章太炎再次给何震写信，云:「接


% 原书第 517 页（PDF 第 530 页）
\BookPage{517}{530}
\noindent 信后已发一函，托购医书各种。闻妹将赴金陵，想近日已在途也。惟自接信至今，已一礼拜，而汇款仍未寄到。家兄既已言明，想亦不致欺诳。但恐由《中外日报》馆转汇，或有支离胶葛之处。妹回沪时，幸细为调查清楚。若仍在《中外日报》馆，不如亲自寄来为安。刘、卞二处消息如何，幸告。兄麟顿首。阳十一月二十五日。」1908年何震注此信，称:「『刘、卞二处』，刘即口口口之姊，系章下狱后，刘允月贴二十金，至今未交者也。卞即前长崎领事卞昌，张之洞之女婿。彼于去岁八月致函张之洞，誓言决不革命，决不与闻政治，且言中国革命决难成功。若赠以巨金，则彼往印度为僧。书与申叔所见，始知彼与官场有往来。及我返国，彼知吾兄何誉生与长崎领事卞昌亲善，彼为张婿，故属我往长崎访之，使再致书于之洞。」

11月30日，《天义》第十一、十二卷合册出版，其「图画」刊载〈世界同胞〉、〈布鲁东像〉、〈饶果像〉、〈鹰山共产村图〉，「论说」刊载申叔〈亚洲现势论〉（附录〈中国现势论〉）、震〈女子非军备主义〉、独应〈论俄国革命与虚无主义之区别〉，「学理」栏刊载申叔〈苦鲁巴特金无政府主义述略〉，「时评」栏刊载扶弱〈防淫奇策〉、独应〈见店头监狱书所感〉、〈中国人之爱国〉，「译丛」栏刊载忱刍〈杜尔斯特致支那人书（节译）〉，「来稿」栏刊载怪汉〈悲佃民〉、〈平民来函〉、〈平原断侵来函〉，「杂记」栏刊载志达〈女界吁天录〉、申叔〈穷民俗谚录〉、申叔〈读书杂记〉，「记事」栏刊载公权〈无政府党大会后事实汇记〉、〈社会主义讲习会记事〉、〈万国革命运动记〉、〈万国女界运动记〉，「附录」栏刊载高亚宾〈废纲论〉、〈革命原理〉（摘录〈新世纪〉、〈家庭与教育〉）。所附〈本社捐金及赞成芳名〉计四人:陆君每月一元，剑侠君十元，陈楚楠君五元，黄礼君五元。

何震在《女子非军备主义》一文中指出，「夫非军备主义所以有益于人民者，盖昔日人民革命，咸依暴动政策，筑塞而守；今则都市街衢广阔，据守甚艰。且政府军队武器，愈加整顿，以保护政府、资本家，非人民所克抗。非取消极主义，解散军队，则民党未易奏功。且军人亦平民之一，乃以服从为性质，甘居奴隶，至为政府、资本家所利用，日与国内外平民为敌，而趋之于死，盖人类之极可悯者，故


% 原书第 518 页（PDF 第 531 页）
\BookPage{518}{531}
\noindent 不惜竭力运动，使之自觉。是则非军备主义所以助平民攫取自由，而因以保人类之身命者也。今世界多数之女子，非居乎民之地位者乎？非俨然人类之一乎？然即此多数之女子询之，果欲获取自由之幸福耶，抑甘于压抑、永沦奴隶之境遇耶？果欲生命之永保耶，抑甘罹祸难，以受非命之惨死耶？盖奴隶之境遇，非命之惨死，均女子之所深恶者也。既知深恶，则非军备主义当为女子所欢迎，彰彰明矣。」「夫军备为女子之害，非仅一端。」「用兵之国，无论胜败，其影响所及，均为妇女之不利。」「且生民之幸福，不外安乐、平和而已。若用兵之国，为女子者。」「其惨况不可胜言:一曰生离之苦，二曰久别之苦，三曰战死闻耗之苦。」「国家用兵之际，死别生离之惨，毕集于女子之一身，非惟为爱情之公敌也，且将陷女子于贫困死亡之一境。即使出师以后，奏凯生还，然此胜则彼败，战败之国，死者必众。盖死者既为他国之平民，则困苦者，亦为他国平民之家室。吾辈既以救济同胞为目的，岂忍寡他国人民之妻，孤他国人民之子女，以逞国家虚伪之光荣哉？盖战争一事，与妇女均有直接之不利，乃女子所首当反对者也。」所以，「兵备一日不废，吾恐女子因贫穷而陷身危苦者，亦永无解脱之望。」

何震认为，「自古及今，男女之权利所以不能平等者，则以男服兵役，女子独否，故尊卑阶级，亦援是而区。且上古之男子，所以夺掠女子，视为俘虏，而行多妻之制者，亦以用兵之故。盖男子富于勇健之精神，其强力足以制妇女，故为女子者屈服于男子之下。积时既久，遂以服从为性质，而仰男子之指挥。此实男女不平等之原因也。故男女不平等之制，亦由用兵而生。且今之禁遏女权者，犹执女子不服兵之说，以证女子权利，不能骤与男相平。是则男尊女卑之观念，亦由服兵而后生。然男子服兵，既历数千年之久，女子欲分任其责，必非旦夕所能平。若实行非军备主义，使男子解服兵之役，退与女平，既不能执保国之空名，以傲女界，复不能逞其强力，而强女子以服从。此实男女平等之权舆，亦女子脱男子羁绊之初桄也。故为女子者，不欲言男女平等则已，如欲实行男女平等，则舍非兵备主义外，决无运动之方。此固不易之理也。」

因此，何震认为:「非军备主义者，弱种平民女子之大利也。非军备主义行，则弱种泯强种之侵凌，平民脱国家之压制，为女子者，亦可脱男子之羁绊，以博自


% 原书第 519 页（PDF 第 532 页）
\BookPage{519}{532}
\noindent 由之幸福：此实世界和平安乐之先声也。」

12月4日，苏曼殊在上海致信刘三，称：「顷须俟剑妹（即何震）来，方能择日东行。剑妹十五日回乡，云一周可返。今逾半月尚未来，殊悒悒。」又称:「申叔有意明春返居沪渎，以留东费用繁浩，且其老太太适异国，诸凡不便故也。」

12月上旬，刘师培由日本返上海，何震也由扬州返上海，与苏曼殊、柳亚子、朱少屏、韩少屏、邓秋枚、林立山等八人相聚，有合影（载柳亚子编、北新书局出版《苏曼殊全集》）。其间，何震运动张之洞未果，转而利用刘氏数世与两江总督幕府的关系，与端方联系，为章太炎筹集赴印度经费。

中旬，何震、刘师培与柳亚子等在上海张园聚会，酝酿发起南社，刘师培有〈张园〉诗。柳亚子〈南社纪略〉说：「南社的人物，除掉后来作为南社发起人的陈巢南、高天梅和我，次第加入社籍的黄晦闻、朱少屏、沈道非、张聘斋以外，还有刘申叔、何志剑、杨笃生、邓秋枚四人，笃生和秋枚后来始终没有加入社籍。」

12月29日，章太炎收到何震、刘师培夫妇自上海寄去的关于运动赴印度研究佛经路费的信。次日，章太炎回信给刘师培、何震，云：「四、六君鉴：二十九日接得手示，知四弟在船甚苦。黄海风浪，素自平静。今忽狂暴，亦应数中所言也（黄杨厄闰苦虔刘）。东坡已归长崎，亦知二君意中，恐其妄行泄事耳，亦未直告。讲习会开时，大杉来会，讲「巴枯宁」（此三字原文用日本假名一何震注）联邦主义，甚有理致。四弟既不往宁，在沪交涉亦善。前书言恐有枝节，愚意可密致杨仁山书，令其转圜。得则为高阳酒徒，不得亦市南宜僚，不泄人言以求媚也。发愤扬灵，最宜深戒。此行亦由一祸。四、六二君，无可说也。勉为其难，吾所属望。家款尚未汇到，则由误交钱庄之故。此间本月开销尚足，勿念。手此即问近好。毛一顿首。十二月三十日。」何震后来解释说:「四君、四弟，均申叔也。六君，即何震也。毛一，章自称也。东坡者，何震英文师苏曼殊也。刘申叔至沪时，恐章报告清政府，加以暗害，遂伪致书于彼，言已在沪为为彼运动。故章书有『在沪交涉亦

\vfill
\begin{BookNotes}
\BookNote{7.}{刘师培曾祖父文淇、祖父毓崧、伯父寿曾均参与金陵书局校书工作，与两江总督府渊源较深。1902年10月南京省试后，刘师培滞留南京，即于25日与高尊士、徐乃昌拜访缪荃孙，次日缪荃孙又回访刘师培。1903年刘师培会试落第后，于5月21日又与梁慕韩、龚艾堂在南京拜访缪荃孙，并载《艺风老人日记》中。}
\end{BookNotes}

% 原书第 520 页（PDF 第 533 页）
\BookPage{520}{533}
\noindent 善』语。又恐彼以必成相要，故以『恐有枝节』告。杨仁山者，池州杨文会也，以通佛学闻，南京官场多敬之。故彼欲嘱何致书于彼，请其向江督为彼乞恩。由『愚意』至『深戒』，均言此事。又视『此行亦由一祸』以下数语，则彼以运动官场事谆谆相托，不烦言而解矣。」

12月30日，《天义》第十三、十四卷合刊出版，其「图画」栏载〈俄女杰方锦像〉、〈强者受刑之图〉之一、之二、〈马场路自由民乐园图〉，「社说」栏刊载震〈女子教育问题〉、震（目录题作「志达」）〈女子革命与经济革命〉、申叔〈社会主义与议会政策〉、续载〈中国民生问题（市业与野业）〉，「学理」栏续载〈苦鲁巴特金学说述略〉，「时评」栏载志达〈请看日本工女之赁金额〉、〈异哉万国社会党对于女子选举权之提案〉，「译丛」栏刊载苦鲁巴特金着、申叔译〈快愉之劳动〉，「来稿」栏刊载划平王〈世界平等政府谈〉，「杂记」栏刊载自觉生〈女界杂事谈〉(〈美国富室女子之骄愚〉、〈德国女子提倡罢工〉、〈澳州家庭劳动者组合〉、〈俄社会党关于女权之议案〉、〈劳动组合与女子〉），续载申叔〈贫民俗谚录〉，「记事」栏刊载公权〈记事七则〉（〈记美国拟放逐哥尔多曼女史事〉、〈记浦盐斯德俄女子加入革命运动事〉、〈记英国妇人社会党本部〉、〈汇记万国革命风潮〉、〈社会主义讲习会记事〉、〈记河南杂志社拟刊女报事〉），「专件」栏刊载〈神州女报发刊词〉，「附录」栏刊载巢南〈轩亭吊秋文〉、申叔〈从军苦歌〉（七首）、〈滇民逃荒行〉。所附〈本社捐金及赞成芳名〉计列五人:复汉君二元，何君三十元，吴女士廿五元，倚剑君三元，无生君一元

何震在《女子教育问题》中指出:「近日女子教育，均奴隶之教育也。不惟亚东为然，即欧美亦然。故欲谋女子之解放，非惟形式之解放已也，尤当求其思想上之解放。所谓思想上之解放者，即排除一切奴隶教育是也。」

「考现今东西各国之女校，其教科首崇伦理，而西洋或代以宗教。夫西教之宗旨，固以女子为罪恶之原，谓生人犯罪，由于女子者也。固以女子为服从己夫，如从天主者也。因以男子象神而造，女子则否，当永久服从男子者也。使女子而受宗教教育，势必以卑贱自甘。」「由古及今，女子因身受宗教教育，以致妨遏思想


% 原书第 521 页（PDF 第 534 页）
\BookPage{521}{534}
\noindent 者，不知凡几。故近日意国女子，于宗教教育，已要求废止，诚以宗教教育，不啻陷女子于奴隶。此女子教育首当改良者一也。」

「现今女子教育权，握于男子之手，即一切教科，亦男子制定。故男子仅为便利一己计，而不复计及女子之幸福。例如宗教、伦理二科，所以阻碍女子之思想，而使之不敢抗抵男子者也。家政一科，欲使其治理家庭，整然有序，而省男子之劳也。工艺一科，欲使之俯就赁金制度，不倚男子之养赡，而男子得省其内顾之忧也。其目的所在，夫固路人皆知矣。于此而犹曰助女子兴教育，其谁信之？」「今为女子解放计，必屏斥一切之旧教育，而代以新教育。盖宗教、家政二科，必宜废弃者也；伦理一科，亦宜屏除往昔思想，而代以新思想，即排斥昔日腐败道德，而导以互助、博爱新道德是也。若工艺一科，宜更易其学习此科之目的。此则吾辈对于教育问题之宗旨也。」

「至于历史各科，均为启发思想之用。惟关于国家主义者，宜在屏遗之列。理化、博物各科，尤为重要。若音乐、图画，本女子之特长，然以待役于人及娱悦男子为目的，则与卖艺无异。此宜导以高尚之思想者也。又中国女子所唱乐歌，多含奴隶主义，亦今日所当改良者。女子肄习医科，亦深合博爱之旨。惟欧美、日本各国，凡女子肄习此科者，多给役军旅之中，以为看护妇，此则深可嫉视者也。政治、陆军、员警各学科，均专属于男子者也。今为消灭人治、排斥军备计，则女子欲与男子平等，不在争习此科，惟有实行无政府、非军备之运动耳。惟农学一端，关于民生者甚巨，乃女子所当从事者也（若商业科，亦不必肄习）。近法国女子多争习农科，此世界女子所宜效法者。依此而行，庶一般妇女，不为奴隶教育所囿，而导其自觉之心。盖昔日之教育，均造成奴隶者也，故女子奴隶教育，与女子奴隶制度相伴而生。嗣今而降，新社会之制，亦与新教育制度相伴而生，乃以造成自由人民为目的者也。既以造成自由人民为目的，以实行家庭、国家、社会大革命，则教育革命，实一切革命之权舆。凡在女界，可不从事于斯乎？」

何震认为：「幼儿当为共有，非家族所得私。幼稚教育，当属之社会，不当属之女子。兴女学之目的，所以为女子解放计，非谓学成以后，当幽闭家庭，而为男子尽教子之义务也。故中国所兴各女学校，其宗旨目的，竟无一而非误。又何怪其


% 原书第 522 页（PDF 第 535 页）
\BookPage{522}{535}
\noindent 野蛮章程订于政府之手乎？」何震的结论是:「家庭革命为一切革命之基。欲兴家庭革命，必自兴女学始。」

1908年二十三岁

1月12日，何震与刘师培、柳亚子、杨笃生、邓实、黄节、陈去病、高旭、朱少屏、沈砺、张家珍等于上海酒楼小饮，相约结社并摄影。柳亚子有〈偕刘申叔、何志剑、杨笃生、邓秋枚、黄晦闻、陈巢南、高天梅、朱少屏、沈道非、张聘斋海上酒楼小饮，约为结社之举，即席赋此〉:「慷慨苏菲亚，艰难布鲁东，佳人真绝世，余子亦英雄。忧患平生事，文章感慨中。相逢拼以醉，莫族酒樽空。」高旭作〈丁未十二月九日国光雅集写真，题两绝句〉云:「伤心几复风流尽，忽忽于兹三百年。记取岁寒松柏操，后贤岂必逊前贤。」「余子文章成画饼，习斋学派断堪师。荒江岁暮犹相见，衰柳残阳又一时。」（《神州日报》1908年3月7日)

1月15日，《天义》第十五卷出版，其「图画」栏刊载〈伏诛各元首之肖像〉、〈波兰革命图〉，「论说」栏刊载震〈论中国女子所受之惨毒〉，续载〈社会主义与国会政策〉，「学理」栏刊载恩格尔斯作、民鸣译〈共产党宣言序言〉，「时评」栏刊载志达〈秋瑾死役之冤〉，「译丛」栏刊载申叔译〈未来社会之生产方法及手段〉，「来稿」栏刊载〈山东千佛山奇狱记〉、〈金华织布局虐待工女记〉，「杂记」栏续载〈女界吁天录〉、〈贫民俗谚录〉。

何震在《论中国女子所受之惨毒》中指出，「吾辈之所主张者，扩张女权也，恶男子以强权加之女子也。夫女子为男子强权所加，固为可闵。若夫己为女子，不能抗抵男权，而徒以横暴之强权，加于服属己身之女子，则其惨毒之罪，尤属可诛。故吾辈之旨，不惟排斥男子对于女子所施之强权，并反抗女子对于女子所施之强权。」

何震认为:「夫今日中国一般之妇女，所受惨毒，为五洲万国所未闻，则以女子之中，有富有贫，有尊有卑。贫者受制于富，卑者受制于尊，而富者之于贫，尊者之于卑，待遇之酷，遂暗无天日，使颠连无告之女子，不惟受制于男，亦且受制于女。则谓此境为世界最惨之社会可也，则谓此苦为中国女子特殊之苦亦可也。」


% 原书第 523 页（PDF 第 536 页）
\BookPage{523}{536}
\noindent 何震还详细列举了富室虐待婢女的七种原因和常见的十五种非人严刑，以及「童养媳之制」形成的四种原因和君姑责罚童养媳的七种恶行。

1月中旬，章太炎第四次写信给还在国内的何震刘师培夫妇，叮嘱他们注意保密。信中说:「四、六弟如见:接信具悉一切。家款近已汇到，家中亦皆安堵。六弟为我尽力，切至周详，感甚。黄、叶亦无他语，惜已明知四弟到沪，在外喧传。黄更知兄欲出家，则此事已稍漏泄矣。运动之事，想二子无不周知。口口口、口口口由沪来东，乃欲为《新世纪》设发行所，别无他故。此报语多模糊。盖留学欧洲者，不知中国情形，辄以欧美相拟，亦犹留学日本者，不知中国情形，辄以日本相拟耳。凡议论不在高深，只求划切。高而不切，言无政府，与言十六字心传无异。兄已详告周君，并望四弟转语笃公也。手此即颂近祉！兄毛一顿首。」何震曾为此信加注云:「彼接何震信，不知已受何欺，转感其尽力。黄者，湖北人黄抱香也。叶为黄友，浙江人，忘其名。黄见章不做《民报》，将疑彼无心革命，在东京对人宣言，故彼畏之甚，疑其尽知彼事。又叶为邻人，知申叔返沪，故信中言『运动之事，想二子无不周知』也。」接此信后，刘师培作〈上端方书〉，献「弭乱十策」。端方表示，同意刘师培所提出的「按月支给」章太炎款。刘师培即写信给章太炎，称卞商端方，同意「按月支付」。

30日，章太炎接到刘师培信后当即复信，反对「按月支给」办法。称：「领事按月支款之说，万难允从。一、若按年分摊，则一岁不过千余圆，或仅数百，必不敷用。二、若摊年过久，章甫去江宁后，事即中寝。三、领事为政府所派，非两江私派。若果迁延抵赖，亦无如何。以留学官费证之可见。要之，不以意气相斯，尽力磋磨，亦无益也。弟若转圆，当要以先付三分之二，不则二分取一。如或未能，当面回复。此则当令六弟任之。手此敬问起居不具。兄毛一白。阳三十日。」何震注云:「此信最有关系。何接彼第四信，复致书于彼，故为疑问之词，谓『将付领巨款欤，抑至印按月支款欤？请明示，以便同交涉』。彼乃以此函相答，反对『按月支款』之说。章甫者，端方也。」

此次交涉因双方意见不一，即告失败，刘师培夫妇则入端方之彀，充当暗探，破坏革命。冯自由〈记刘光汉变节始末〉云:「光汉心衔所提议改组同盟会、攘夺


% 原书第 524 页（PDF 第 537 页）
\BookPage{524}{537}
\noindent 干事职权之策不成，渐有异志。会是时清吏肃王善耆、铁良、端方等鉴于徐锡麟之暗杀，咸有戒心。各自设法向党人施展金钱政策，使为己用。肃王首罗致皖省党人程家柽，遣至东京，愿贡献万金于同盟会本部，以尽友谊，声明决无条件。刘揆等以渴不饮盗泉拒之。端方亦授意败类学生多人从中作祟，或伺隙离间，或用金钱收买。党人意志薄弱者，间为所愚:光汉之妇何震及姻弟汪公权即其人也。何、汪不独从此入于侦探一流，且形同夫妇，宣言公夫公妻不讳。适光汉因事与章太炎、陶成章等大起冲突，又以勾结日人谋夺党权事不理人口，何、汪等乘之，日夜惩恿光汉，使入官场，以图报复。光汉外恨党人，内惧艳妻，遂不得不铤而走险，始真为江督端方之侦探矣。」8案，冯自由将刘师培变节原因归咎于何震、汪公权的恿，根据不足。刘师培夫妇回国期间，汪公权一直在日本。1907年12月的社会主义讲习会两次集会，汪都参加并作了报告。其次，刘氏变节，也不能仅仅归咎于何震。章太炎托刘师培夫妇运动赴印度学佛路费，为刘氏变节提供了机会。再次，此时刘氏夫妇也因为留日费用浩大，加上办报、开会之需，经费拮据，因而此番回国筹款，也为决定是否住在日本作一选择。他们选中端方作为目标，可能是端方在朝中得势，在预备立宪中又呈上了梁启超代拟的二十万言奏稿，看上去还是一个开明官吏。当然，端方在这笔政治交易中给了刘氏夫妇一笔钱，使得他们能在日本又住了将近一年。另一方面，刘师培与同盟会党人的政治分歧已经很深，早就有了「维新不如守旧」的念头，政治上的分道扬镳是必然要发生的。此外，冯自由因反对共产党、反对国共合作，10因此对社会主义、共产主义理论一直持反对态度。他所谓的「形同夫妇，宣言公夫公妻不讳」，均为悬测之词，不足为据。

2月中旬，何震、刘师培返回日本。26日，苏曼殊由日本寄信给刘三，云：「申叔伉俪西来，询知足下平安。前此寄上一信，并日文杂志四册，《天义》第

\vfill
\begin{BookNotes}
\BookNote{8.}{冯自由，《革命逸史》(上)，页332-333。}
\BookNote{9.}{南桂馨，〈山西辛亥革命前后的回忆〉称:「刘师培全家在日，经常闹穷。记得有一次，他向《民报》社借三十元，宋教仁以无款可付，未能照办。章炳麟出而打抱不平，和宋争执得很激烈，刘揆一也很为难。经我和张继从中调处，才算了结。」}
\BookNote{10.}{冯自由反对共产主义、反对共产党的观点，在其所着《革命逸史》的〈章太炎等之护党救国公函〉、〈秦毓鎏事略〉等文章中均有所反映。直到20世纪30年代，一些反对共产主义的人，其强加给共产主义的罪名中，仍然有「共产共妻」这一条。}
\end{BookNotes}

% 原书第 525 页（PDF 第 538 页）
\BookPage{525}{538}
\noindent 八、九、十合一册，均托秋枚转致，妥收未？曼现暂寓神田猿乐町一丁目二番地清寿馆。日间须往横滨病院静养，盖肝跳症也。足下赐教乞寄申叔转交无误。申叔下月迁居，曼病愈后亦同往。申叔云:足下会春尚至杭，不致多生烦恼矣。」

3月20日，社会主义讲习会改名「齐民社」，并召开第一次会议。

3月底，《天义》第十六至十九卷合册（春季增刊号）出版，其「图画」栏刊载〈实行暗杀九义士之肖像〉（即罗石、濮鲁西、克舍利、柏涛、查古斯、卢支尼、赵沙滨、路佛、哥萨多）、〈葡国王后太子遇刺图〉、〈石门华氏肖像〉、〈曼殊画四幅〉，其「译书」栏刊载马尔克斯、因尔格斯作、民鸣译〈共产党宣言〉（前有申叔〈序〉，后有申叔案语），哈因秃曼作、齐民社仝人译〈社会主义经济论〉，布利斯作〈社会主义史大纲〉，苦鲁巴金作、申叔译〈面包略夺〉，齐民社同人译〈无政府主义之哲理同理想〉，杜尔斯德作〈俄国革命之趣旨〉，杜尔斯德作、忱刍译《致中国人书》，怪汉译〈俄国第二议会提议之土地各法案〉（前有申叔《序》）、公权译〈无政府党大会决议案〉、马拉迭斯丹作、雠译〈无政府共产主义之工人问答〉（前有刘申叔〈序〉），「附印」刊载梅毓〈区田考〉（前有刘申叔〈序〉）、「编纂」栏刊载申叔辑〈选举罪恶史〉、志达〈女子问题研究〉、〈贫民唱歌集〉（申叔〈工女怨二章〉、民鸣〈农民哀六章〉、石门华作、申叔译《希望诗二章》，孙枝蔚《善哉行》、谈献《小车行》、无名氏〈述所见〉、郑燮〈逃荒行〉、秋鸡哈威作、张继译〈咏同盟罢工〉）、申叔<Esperanto词例通释〉（前有〈总序〉），「记事」栏刊载〈记葡王被刺事〉、〈万国革命记〉，「附录」栏刊载〈英译吊古战场文〉、〈英译《礼记》孔子过泰山侧一节〉、〈英译汉诗〉（王昌龄〈闺怨〉）。

本期所载〈本社重要广告〉称:「本报从今岁起，每季发行增刊号一巨册，本号所列各译稿，定于夏、秋、冬三季内出全。此外西文各书之译稿，凡此册所未载者，亦于下三号补登。」然《天义》此后即停刊。

4月6日，世界语讲习会在刘师培住宅首次开班。参加者有何震、刘师培、张继、苏曼殊、景梅九及中国留学生20多人，拟每周月、水、金曜日下午四时至五时半上课。


% 原书第 526 页（PDF 第 539 页）
\BookPage{526}{539}
此后，何震、刘师培与章太炎交恶。关于章刘交恶的起因，大致有数说。

一为生活起居琐事说。梅〈青溪旧屋仪征刘氏五世小记〉称:「戊申年又赴日本，这一次外祖母同去的。章太炎也在日，住在舅氏家中。章先生囚首垢面，衣巾经月不浣，养松鼠于袖中。果壳和干肉狼籍，室内虫蝇飞扰。又夜斥下女，继以号哭，旋复高歌。舅母不能堪噪，逐之。先生逡巡出，入小旅舍。外祖母以为同在客中，宜相照顾，晓舅母以大义，命舅氏迎之归。故先生见外祖母，必致敬礼。与舅氏亲若昆季。乃未久，以论学及政见不同，闻其中有奸人播弄，遂略有龃龉。」冯自由〈章太炎事略〉称:「盖章生平不脱名士风尚，视沐浴为畏途，幽囚后，西狱吏每日强之澡浴，故体魄因而日健也。」又称：「章对于自用衣袜，亦常杂乱善忘。」11

二为争学术声名说。梅铉〈青溪旧屋仪征刘氏五世小记〉载:「舅氏与余杭章太炎氏学术既同，情好甚敦。因有不良分子造为诽谤，说『章太炎外和内忌，游扬当道，有不利于孺子之心』等语，登诸报章。再经舅母加以饰词，舅氏引为大恨，遂向太炎绝交。太炎百计修好，舅氏都是置之不理，有信亦不覆的。」章士钊柳文指要·刘师培论古文〉也说:「太炎尝论凌次仲，谓『次仲号为工骈体，然止似文章家之骈体耳，无以自见所学也，彼岂谓一为文章，便当抛却经师本色耶？』此殆欲以晏文献意中之柳子厚，策励申叔。」又称：「申叔吾友也，平日学优于吾，吾甚敬之，独少时违难东京，太炎与吾及申叔聚处时，偶谈文事，意辄相左。时吾撰《中国文典》，太炎谓是别通故训之一途径，而申叔反对甚烈。本年6月1日，章太炎致书孙诒让，也是持这一说法，称:「仪征刘生（旧名师培，新名光汉，字申叔，即恭甫先生从子），江淮之令，素治古文《春秋》，与麟同术，情好无间，独苦年少气盛，喜受浸润之谮。自今三月后，谗人交构，莫能自主，时吐谣该，弃好崇仇。一二交游为之讲解，终勿能济（以学术素不逮刘生故）。先生于彼则父执也，幸被一函，劝其弗争意气，勉治经术，以启后生，与麟戮力支持残局。度刘生必能如命。怀懊陈述，非为一身毁誉之故，独念先汉故言，不绝如线。非有同好，谁与共济？故敢尽其鄙陋，以浼先生。惟先生少留意焉。」

\vfill
\begin{BookNotes}
\BookNote{11.}{冯自由，《革命逸史》（上)，页51、52。}
\end{BookNotes}

% 原书第 527 页（PDF 第 540 页）
\BookPage{527}{540}
三为何汪挑唆说。汪东〈同盟会和《民报》片断回忆〉中说:「一九0八年，同盟会内部曾发生过向南北洋告密的事。不久刘申叔就投靠端方，做了端的幕僚，月薪三百两，因此证明是刘所为。实则其中还有内幕。刘是一个学者，终日埋头著作，又有肺病。他与妻子何震本是表兄妹，而汪公权又是刘的表弟，在日本时同住一起。何既好名，而又多欲。他一面利用刘能写文章，替她出名办了一种月刊（是否即《天义报》记不清），提倡女权，并带无政府色彩；12一面又对刘不满足，行为放荡。汪趁此勾引，与何发生了关系。那时章太炎从民报社迁居刘处，于无意中发觉了汪何的秘密，便私下告诉刘，刘的母亲也听见了，非但不信，反大骂章造谣，离间人家骨肉。这事声张开来，汪、何二人当然恨之切齿，尤其是汪。我那天正去看太炎先生，被刘母遮住了诉一顿冤（那时我尚未做太炎学生）。汪公权满脸凶气，眼睛里都是红丝，跳出跳进，嚷着『我们白刀子进去，红刀子出来』，我简直不懂他说的什么。恰巧刘揆一也在那里（刘可能是去调解这件事的），便同我一道出来，摇头冷笑道:『这种江湖上的下流口吻，他（指汪公权）拿来吓谁？』我是觉得于我无关，不愿再问。随后太炎搬回民报社。」

从章、刘关系看，刘师培、何震因帮助章太炎联络端方而被端方收买，其中固然有刘师培、何震自己立场不坚定的因素，然而，章太炎托其运动赴印度学佛路费，确实是其中原因之一；此事未能成功，章太炎有怨言，也在情理之中。但导致二人决裂的直接原因，似乎仍是学术名声之争。一是刘富曾〈亡侄师培墓志铭》称刘师培「少年气盛，思欲有以自见」，应当是对刘师培个性的中肯评价。二是章太炎在《国粹学报》第41期发表〈驳中国用万国新语说〉，对刘师培、何震正在热心推广的世界语公开表示反对。刘师培即在该刊第47期发表〈论中土文字有益于世界〉，针锋相对地提出:「今欲扩中土文字之用，莫若取《说文》一书，译以Esperanto（即中国人所谓世界语）之文，其译述之例，则首列篆文之形，或并列古文、籀文二体，切以Esperanto之音，拟以Esperanto相当之义，并用彼之文，详加解释，使世界人民均克援中土篆、籀之文，穷其造字之形义，考社会之起源，此亦世界学术进步之一端也。」三是章太炎发表在《民报》第17号〈国家论〉，称:「余向者于社会主义讲习会中，有遮拨国家之论，非徒为期望无政府者说。虽

\vfill
\begin{BookNotes}
\BookNote{12.}{全国政协文史资料研究委员会编，《辛亥革命回忆录》第五集，中华书局，1963年。}
\end{BookNotes}

% 原书第 528 页（PDF 第 541 页）
\BookPage{528}{541}
\noindent 期望有政府者，亦不得不从斯义。然世人多守一隅，以余语为非技过甚。故次录前论，附以后义，令学者得中道观云。」四是章太炎在《民报》第17号发表〈排满平议〉一文，认为「诚欲普度众生，仅一切得平等自由者，言无政府主义，不如言无生主义也。转而向下，为中国应急之方，言无政府主义，不如言民族主义也」。同时，对刘师培一贯作为立论基础的「汉族西来说」表示怀疑，称:「汉族自西方来，非有历史成证，徒以考索比拟而得之。刘师培后来在《衡报》第2号发表〈社会革命与排满革命〉，对此加以反驳。五是本年5月9日，陈去病致书高旭（收入《国学丛选》第六集），邀请高旭、柳亚子到杭州张煌言墓前一哭，兼吊秋瑾，并希望以章太炎、刘师培反目为鉴，永固贞盟，称:「刘、章结果，弟数年前早料到。盖两人皆经生，铿铿好辩，不肯服输，匪若吾侪终不忍以意气坏大局也。好兄弟而如此，能无泫然？所原吾侪当日凛凛，无蹈此覆辙，而自破其贞盟也，则吾心安矣。」同样也是持争学术声名说。六是1909年章太炎致信刘师培希望其迷途知反，称：「仆之于君，艺术素同，气臭相及。猥以形寿有逾，恒人视之，若先一饭。精义冥思，亦有多算。君雅好闻望，不台于先我。自谓文学绪业，两无独胜。怀此觖望，弥以恨恨。然仆岂有雍蔽之志哉!学业步骤，与年相将。悠悠之誉，又非由己。畴昔坐谈，盖尝勤攻君过；时有神悟，则推心归美。此盖朋友善道之常，而君岂忘之耶？」也是强调二人学术名声上的问题。七是南桂馨在〈辛亥革命前后的回忆〉13中也指出:「章主张单纯民族革命，彼以为世界革命远而无望，反而松懈了我们的民族革命事业。刘和我以为联络世界革命党人，对于我们的革命事业有益无害。章大不以为然。彼此的分歧就此发生了。」「刘和章炳也互相轧轹：刘大评章讲小学之无当，章亦反讥。」也认为是因政治主张和学术纷争导致二人反目。八是1909年8月21日蔡元培给吴稚晖的信中，在讨论刘师培变节问题时，说刘师培有好胜、多疑、好用权术三个特点，并举例说:「有一日，吴彦复言:夏穗卿到彦复处，见申叔所作〈攘书〉，有言黎民即汉人指目色苗种之名，则大诧曰:光汉小子，好盗人书。盖穗卿曾于〈汉族纪口〉（见《新民报》）中有此说，以为申叔袭之而不着其所自出也。时陈镜泉在座，曰：申叔前见屠敬山之〈历史讲义〉有此说，尝曰:吾书不可不速刊。否则人将以此说为创于屠氏也。」亦为章、刘学术名

\vfill
\begin{BookNotes}
\BookNote{13.}{全国政协文史资料研究委员会编，《辛亥革命回忆录》第五集，中华书局，1963年。}
\end{BookNotes}

% 原书第 529 页（PDF 第 542 页）
\BookPage{529}{542}
\noindent 声之争而发。

4月21日，何震致书吴稚晖，揭发章太炎与端方关系。函称:「稚晖先生大鉴：久慕大名，恨未晤面，以聆教诲，怅甚！留法中人发起《新世纪》，久为敝等所崇拜。不意日京民报社之章炳麟，因与先生有隙，即加诋毁；又因恨《新世纪》之故，并憾及留法之人，以及法国各学派。于此次张君溥泉如法，彼即虚造伪言，甚至欲将渠入狱，一月引渡清国，险恶如此，罪当如何！而东方无知之革命党受其影响，亦排斥无政府主义及世界语，故贵报于东方不克发达，而观19期《民报》，载有答先生书一篇，痛加宣布。其言虚实，鄙人固不能知，但彼暧昧之历史，则知之甚晰。试陈之以备参考。章炳麟，一名绛，字太炎，又字枚叔，别号末底、西狩、载角，浙江余杭人。幼婴羊疯疾（今尚缺二门牙），甫应县试，其疾大作，遂纳粟为国子生，且从伯兄习制艺，冀应乡举。则其革命思想，非具于壮年之前，且非蓄排满主义，始以应试为耻，彰彰明矣（此事彼家族及炳麟亲对吾言）!彼又受张之洞之招，供其役使。又皖人吴保初为故提督吴长庆子，彼在沪常主其家，则又非疾视官场者比。且彼庚子年偕保皇党上书李鸿章（此书由章起稿），又致书张之洞及江南道员俞明震，多以变法冀清廷（今此稿犹存），并明震复言，将此面呈老帅，老帅大悦云（即刘坤一）。去岁曾受铁良二百金（系由国事侦探程家柽经手，刘林生言），又去年九月上张之洞书，与伸旧谊，逢迎其国学，末言若助以巨金，则彼于政治问题，不复闻问，并谢辞《民报》编辑（此言系下婢名□□所发）。余甚多，不克枚举。近已用针笔板照像法付印，俟成即寄上。外附《民报》19期答先生原函寄上。是端已为同志中一友人披过。如合尊意，留登贵报。余容续布。即颂自由幸福并祝《新世纪》无政府万岁！Chin上。西4月21日。」吴稚晖曾将其章、刘反目情形函告蔡元培。蔡元培复信说:「枚叔末路如此，可叹可怜！然申叔亦太不留余地。公所谓贻反对党骇笑，诚不免。公及博公劝之，甚善。」14

4月28日，刘师培在东京出版《衡报》第一号（创刊号），为躲避日本警方的干涉，托名澳门出版。

5月上旬，受刘师培与章太炎关系破裂的牵连，本与刘师培夫妇同住的苏曼殊

\vfill
\begin{BookNotes}
\BookNote{14.}{孙常炜编，《蔡元培先生全集·遗墨》，台湾商务印书馆，1991年。}
\end{BookNotes}

% 原书第 530 页（PDF 第 543 页）
\BookPage{530}{543}
\noindent 迁居友人处避祸；章太炎即搬入牛込区新小川町二丁目八番地民报社内居住。7日，苏曼殊在日本写信给刘三，称:「盖近日心绪乱甚，太、少两公（指章太炎、刘师培）又有龃龉之事，而少公举家迁怒于余。余现已迁出，飘泊无以为计，欲返粤一转，奈无资斧何！故只可沿门托钵。」

5月24日，章太炎在《神州日报》发表〈炳麟启事〉，称:「世风卑靡，营利竞巧。立宪、革命，两难成就。遗弃世事，不撄尘网，固夙志所存也。近有假鄙名登报或结合者，均是子虚。嗣后闭门却扫，研精释典。本日即延高僧剃度，超出凡尘。无论新故诸友，如以此事见问者，概行谢绝。特此昭告，并希谅察。章炳麟。」

6月1日，章太炎致书孙诒让，托其调停与刘师培的矛盾。18日，汪公权由日本起程赴爪哇岛。10日，章太炎发表〈特别广告〉，否认〈炳麟启事〉出自其手笔。后人以为〈炳麟启事〉为刘师培夫妇伪作。

6月中旬，何震由日本回国，筹措《衡报》办理发行手续的保证金。

8月8日，《衡报》第10号刊载〈何震女史由镇江来函〉，云:「由上海至镇江，所经沪宁铁路，凡车站之中，不准卖食者入内，卖者均于站外高呼，然下车购物者有几？故此等之小贩，生活极为艰难，而车中食物，则贵于市间数倍。一盂红茶，亦索价五分，则贫民之乘车者，亦惟忍饥忍渴而已。又各小站仅一面有车台，下车者必经轨道，稍不注意，即为车轹死，闻毙者已数人矣。扬、镇一带，近日物价愈贵。扬州非商埠，然鸡蛋一枚，已贵至十文以上；猪肉一斤，则需值一百六十文。较之前三年，殆增一倍，则贫民之多，可知矣。镇江各商店闭歇者甚多。其原因均因钱庄盘剥，负债不克偿，以致歇业。然歇业之后，每店之中，均有十余人赋闲。闻前日镇江一带破圩甚多，现遍处皆饥民矣。」

8月21日，日本曲町员警署传唤刘师培，认为《衡报》发行手续不全，要求正式发行前需补齐有关手续。刘师培要求给予10天的期限。

8月31日，日本曲町员警署再次传唤刘师培，要求9月8日前交纳保证金，方可正式发行《衡报》。


% 原书第 531 页（PDF 第 544 页）
\BookPage{531}{544}
9月6日，何震自国内回日本，所筹保证金数量不足。

9月8日，刘师培夫妇向日本曲町员警署交纳筹集到的《衡报》发行保证金，但仅为员警署要求的一半，不得不继续奔走筹集。

9月9日，日本曲町警方跟踪调查刘师培、何震筹集《衡报》发行保证金的动向，并注意到了《衡报》托名澳门出版的问题。本日，日本警方的调查资料〈劉光漢行動二就テ〉1称:「刘光汉于研究各国社会主义者著作外，且承担澳门平民社所发行之《衡报》通信事务。警方因对《衡报》发行处略有怀疑，故作调查，发现该刊发行全在刘氏住所。上月二十一日，所辖曲町员警署对其传讯，并以其缺少发行手续为由进行训诫。彼因办理正式手续继续发行抑或就此停刊，尚需认真考虑，乃请求准予十天考虑时间。警署准其请，令其于上月三十一日至本月八日间交付保证金，并经审查同意后，方可正式发行。其妻何震如前所报六月中旬秘密回国，三日前回东京，刘氏拟待其妻回京，以其所筹之款支付保证金。然其妻所携之款仅足其半，现正多方奔走筹款，尚无进展。」

9月15日，刘师培到曲町员警署，补交齐所需的《衡报》发行保证金，获得正式发行手续。但日本警方继续监视其活动。本日，日本警方的调查资料《劉光漢二就テ》16称:「前报清国革命党人刘光汉本日至本厅，为《衡报》正式发行交纳保证金。」

10月10日，《衡报》被日本警方查封。

10月19日，《民报》也被日本警方查封。

10月20日，刘师培、何震开始作回国的准备。日本警方调查资料〈劉光漢言動二就テ〉17称:「清国革命党员、社会主义者刘光汉，以其所营之《衡报》第一号于十日前被迫停刊，知其在日本已发展无望。拟远行欧洲或印度，与彼地同志者商议再发行一大机关报，以成其素志。」

\vfill
\begin{BookNotes}
\BookNote{15.}{译自嵯峨隆着，《近代中国革命的幻影—刘师培的思想与生涯》，研文出版，1996年，页155-156。}
\BookNote{16.}{译自嵯峨隆着，《近代中国革命的幻影—刘师培的思想与生涯》，研文出版，1996年，页156。}
\BookNote{17.}{同上，页157。}
\end{BookNotes}

% 原书第 532 页（PDF 第 545 页）
\BookPage{532}{545}
10月31日上午8时，何震赴神户，准备回国，刘师培和母亲李汝仍留在东京市曲町区饭田町6-31号齐民社内。日本警方调查资料〈劉光漢ノ妻歸國ノ件〉18称:「清国革命党人刘光汉（别名刘申叔）之妻何震本次突然回国。本日上午八时自新桥赴神户，刘为其送行至本地。据其至神户后之情形推测，允或同时回国。」

11月2日，何震由神户到达长崎，住广场马町四海楼。日本长崎县知事〈清國革命黨員二關スル件〉1称:「顷接函，称清国革命党人刘光汉外，同国人一男一女，同行日本人一名，昨夜自神户出发，当日下午三时四十分于门司换乘长崎列车。次日下午四时四十五分，佐贺员警署以电话相通报。同日下午八时该员警署又以电话通报，称清国革命党人现已通过佐贺车站。然除刘光汉外，清国男子一名，着和服者二名，同乘二等列车，着玄色护士服清国女子一名，同在二等车厢，与原电话通报所称日本人一名不符。该车于当日下午十时五十一分至长崎，遂作检查。仅有着护士服之清国女子，并无他人。同人称为何志剑（二十三岁），住市内广马场町四海楼。二日后下午五时，乘当地航班春日丸号赴上海。停留中，曾电告东京市曲町区饭田町六一卅一齐民社内刘老夫人，告以即将出发，并无他言。至前称刘光汉者是否到达，经严密检查，至今无到达迹象。二日后，又接山口县查询清国革命党人刘光汉是否到达之电报，经询与清国女子同车之本市本下町碳商诹访安太郎（年约五十），彼谓自大阪登车，于广岛车站附近见车内一着洋装女子，该女子称赴长崎，因无熟人，至长崎已夜深，因问投宿何所。乃语言不通，仅可笔谈。该女子住广马场町四海楼一夜，并无其他同行之人。又门司车站发车前临时检查，警方询彼该女子有无同行者，乃称自广岛以西均与该女子同车，未见同行之人，亦未见与彼交谈者。」

11月4日下午5时，何震乘春日丸号轮船回国。6日下午6时，刘师培母亲李汝从新桥出发回国。

11月9日，日本警视总监〈劉光漢二就テ20报告称:「清国革命党人刘光汉

\vfill
\begin{BookNotes}
\BookNote{18.}{同上，页157-158。}
\BookNote{19.}{同上，页158-160。}
\BookNote{20.}{同上，页160-161。}
\end{BookNotes}

% 原书第 533 页（PDF 第 546 页）
\BookPage{533}{546}
\noindent 因送其妻何震回国，上月三十一日至神户。据至本地情形分析，其妻一旦回国，彼允或同时回国，均已前报。彼送其妻至下关后（其妻自长崎赴上海），曾电告其母，称彼将赴欧洲。刘母七日后取出家财，同日午后六时自新桥赴上海。上月十日《衡报》责令停刊后，彼见在本地发展无望，拟赴欧洲或印度，与当地同志者商计发行一大机关报。若此法可行，则赴欧洲。」

随后，刘师培也返回国内。11月14日，苏曼殊在南京给刘三写信称:「少公（即刘师培）已返国。衲前日过沪，日余即返。闻佩公亦于月杪至沪。」

刘师培回国后，拟在上海开设世界语事务所，传授世界语，并遍告同仁，称：「鄙人因发行《衡报》遭日本政府迫害，拟迁报社至沪，期便传布。至沪后，一俟根基稳固，即秘密出版，并设世界新语事务所，以为通信机关。凡有海内外同志欲通讯、寄报及邮报者，请寄上海租界太古马德星第一千一百五十四号。」21

12月，陈去病自粤返沪，徐自华、刘师培、何震、苏曼殊、包天笑、邓实、朱少屏、沈砺、杨天骥、叶楚伧、林寒碧、徐小淑等置酒相迎，陈去病为述南游状况，不胜感慨。寄尘（徐自华）〈金缕曲·春江即事〉22序云:「巢南先生归自粤中，同人咸集，酒酣，先生为述游迹，慷慨若不胜情。会余亦有所感激，即席塡此，示无畏、志剑、曼殊、天笑、秋枚、屏子、嗲公、千里、楚伧、亮奇及家妹小淑、女儿馨丽。」

这时，刘师培、何震开始充当端方密探，向端方出卖江浙革命党人起义计划，张恭被捕，起义流产。陶成章〈浙案纪略〉说:「清帝死后，光汉意成章回国，日与两江督标军官米占元往各船坞查成章行踪。久之不得，无以复端方之命，而以张恭报告于端方，张恭遂被拿问。」

1909年二十四岁

1月2日，刘师培、何震回扬州小住，不久又返回上海。本日，苏曼殊本日在上

\vfill
\begin{BookNotes}
\BookNote{21.}{同上，页161。}
\BookNote{22.}{见《中华新报》1909年4月1日。}
\end{BookNotes}

% 原书第 534 页（PDF 第 547 页）
\BookPage{534}{547}
\noindent 海写信给刘三，云:「申公伉俪今日归家，数日即来。衲亦束装东归。」

张恭被捕后，上海革命党人非常愤怒，派人查寻叛逆者。不久，王金发查得为刘师培、汪公权所为，即将汪公权暗杀于上海。冯自由《记刘光汉变节始末》23谓：「王金发侦知为光汉所为，怒挟枪访之，责其变节卖友，将处以死刑。光汉跪地乞命，谓必以一己生命，保全张恭。恭因得移禁上元狱，幸不死。光汉由是不敢再至上海。汪公权以为无虑，仍时至上海侦探党人举动，卒为王金发枪毙示儆。闻者快之。案发后，王金发偕其友一人，亡命香港，访冯自由求庇。冯乃匿之于湾仔东海旁街七十六号四楼自宅。汪东〈致黄焯书〉亦云:「刘申叔告密于南北洋。袁世凯置之不理，端方则密侦党人，陈陶遗回国即被逮。刘本书生，告密事皆受其妻何震及表弟汪公权二人所劫持。汪至上海，即被王金发刺死。王盖受党中央命令也。」（汪东〈致黄焯书〉，未刊，承黄焯先生录示）

刘师培劣迹败露后，在上海无法容身，于4月到南京，任两江师范学堂教习，入端方幕，「为方考定金石，称匋斋师，名遂替」24。同时在两江师范学堂任教的还有雷恒、陈三立、王伯沆、柳诒征、钟钟山等。〈刘氏五世小记〉称:「端方为两江总督，李瑞清为两江师范学堂监督。这时要开办历史、地理选科，为全国高等教育的先河，必需请硕学高名的人担任教授。访得上海学通中西的姚文栋先生之子姚明辉，夙承家学，聘为地理教授。惟历史一门，仍乏通才。有人建议延聘舅氏。但李瑞清以舅氏名挂党籍，不敢专主。一日，谒端方于宝华盦。端方因得到海内止存三本的《西岳华山碑》，宝爱非常，遂在署内辟一精舍，名宝华盦，公余常与一班文学名流在盦内读碑谈艺。李瑞清先商之丹阳陈庆年。陈字善余，是一个博学多识的人，与端方契洽。时在督署为首席幕僚，言听计从。他本与舅氏有旧，听了极为赞成，力任进言。次日即与端方谈到:『仪征刘氏，三世传经，家学渊源，为嘉、道以来江淮间第一。他本人又是英年博学。虽为革命党人，近年已不谈种族革命。他若能来，实为上选。』端遂嘱江甯藩司樊云门具函礼聘。由李、陈电约返国。舅氏尚在考虑。这时舅母何震久厌居东，听小人之言，适符她的名利思想，以为能与

\vfill
\begin{BookNotes}
\BookNote{23.}{冯自由，〈记刘光汉变节始末〉，见氏着《革命逸史》（上），页333。}
\BookNote{24.}{汪东，〈刘师培传〉，见《刘申叔先生遗书》卷首。}
\end{BookNotes}

% 原书第 535 页（PDF 第 548 页）
\BookPage{535}{548}
\noindent 官场联系，自然另有出路，遂极力惩恿，加以要胁。舅氏是个疏于世故的人，听她的话，不能坚定立场、权其得失，就贸然返国。」25

6月，端方调任直隶总督兼北洋大臣，报纸公布的随员名单中，列有刘师培之名。刘师培入端方幕开始公开。

7月，刘师培随端方北上天津，任直隶督辕文案、学部咨议官等职，何震随往。

11月，端方因慈禧太后灵柩奉安时派人拍照，触怒隆裕太后，被斥为「恣意任性，不知大体」，并被削职。此后，刘师培、何震寓居天津。

1910年二十五岁

1月22日，《新世纪》遵孙中山嘱，刊载章太炎致刘师培、何震五封信以及何震所加案语。

2月22日，何震生下一女，取名颎，举家欢喜异常。

8月13日，女儿刘颎因病夭折。刘师培、何震十分伤心。刘师培作〈女颎圹铭〉，并作〈伤女颎〉诗二首。此后何震再未生育。

1911年二十六岁

5月，四川保路运动爆发，清政府任命端方为督办粤汉、川汉铁路大臣。

5月22日，刘师培随端方南下，何震随行至汉口。

11月17日，端方被杀，刘师培被资州军政分府拘押。辛亥革命爆发。何震从汉口辗转回北京，生活无着。山西同盟会老人杜羲将何震的窘境告知南桂馨，南桂馨即邀何震赴山西太原，暂住南府，先推荐她到山西女子师范任教，后转任阎锡山的家庭教师。

\vfill
\begin{BookNotes}
\BookNote{25.}{梅着、梅英超整理，《青溪旧屋仪征刘氏五世小记》，上海古籍出版社，2004年，页47-48。}
\end{BookNotes}

% 原书第 536 页（PDF 第 549 页）
\BookPage{536}{549}
12月1日，章太炎在各报发表〈章太炎宣言〉，希望起义军不要杀害刘师培。

12月15日，刘师培到达成都，与吴虞、廖平、谢无量相会。

1912年二十七岁

1月1日，中华民国临时政府在南京成立。

1月11日，章太炎、蔡元培发表〈求刘申叔通信〉，查寻刘师培下落。26日，又电请南京临时政府，要求设法保全刘师培性命。安徽都督府秘书科邓绳侯、陈独秀等也致电南京临时政府，要求「曲为宽宥」。29日，南京临时政府大总统府、教育部分别致电四川都督府，要求派员护送刘师培至南京，「勿苛待」。但刘师培没有返回南京，而是留在四川，任四川国学院院副，讲授《春秋左传》。

春，何震得知刘师培在川消息，由阎锡山、南桂馨各赠川资一百元，自太原出发，南下赴川。何震至成都后，与刘师培相会，并随之居蜀中。

1913年二十八岁

6月，刘师培、何震应南桂馨之邀，准备北上山西。

6月24日，杨南皋夫人宴请何震，为其北上饯行。吴虞夫人曾香祖作陪。吴虞《爱智日记》载:「香祖过杨南皋夫人处便饭，因请有刘申叔夫人何震，故去一谈。」

6月29日，何震、刘师培离开成都，沿江东下。

8月25日，何震、刘师培抵达湖北宜昌。

9月，何震、刘师培抵扬州，小住数日，即赴上海，与在上海的谢无量相会。

10月，何震、刘师培由上海赴山西。

11月，何震、刘师培至山西，住南桂馨家中。后由南桂馨作筏，刘师培任山西


% 原书第 537 页（PDF 第 550 页）
\BookPage{537}{550}
\noindent 都督府顾问，何震仍作阎锡山家庭教师。苏曼殊自日本寄函何震、刘师培，告知在日本情形。

1914年二十九岁

春，因山西都督府改编为将军府，编制缩小，顾问裁撤，刘师培由阎锡山、南桂馨以专电保荐赴北京，请袁世凯任用。袁素闻刘师培之名，而袁的亲信秘书闵尔昌与刘师培既是同乡，又有戚谊，从中揄扬，所以袁命刘师培入京。到北京后，由袁克定带领引觐袁世凯。袁世凯授以公府咨议职。何震随刘师培至北京。

1915年三十岁

8月14日，刘师培与杨度等发起成立筹安会，为袁世凯帝制鼓吹。

10月23日，袁世凯任命刘师培署参政院参政。

11月15日，袁世凯授刘师培上大夫衔。

1916年三十一岁

6月6日，袁世凯在北京病逝。次日，黎元洪接任大总统。29日，黎元洪宣布遵行民国元年的《临时约法》，恢复国会，特任段祺瑞为国务总理。

7月14日，黎元洪下令惩办帝制祸首。李经羲以「爱惜人才」为由，经黎元洪等人首肯，将刘师培、严复列入「宽免」之列。刘师培、何震流寓天津。

12月26日，黎元洪任命蔡元培为北京大学校长。

本年，何震姐姐何采苹在扬州病逝。刘师培作〈何大姑哀赞〉。

1917年三十二岁

秋，何震、刘师培到北京，住老爷庙胡同。刘师培应蔡元培之聘，任北京大


% 原书第 538 页（PDF 第 551 页）
\BookPage{538}{551}
\noindent 学文科教授，兼任文科研究所国文门指导教师。1917-1918学年担任的课程有：中国文学（一年级、二年级，每周各三小时），中国古代文学史（二年级，每周三小时）。所指导的研究科目为「文」、「文学史」，每月第二周和第四周的星期四分别与「文学史」、「文」两个方向的研究员会面1小时（后具体日期有所调整）。后又兼任女子高等师范学校讲师。

从这时开始，刘师培肺病日益加深，到处求医问药，所得薪资大部分用于治病，因而经济上日益困难，有时不得不靠借贷度日。

1919年三十四岁

11月20日，刘师培在北京病逝。据《申报》11月27日〈京学界要人之凋谢》称：「身后萧条，一无所有。历二日，始入殓。所有殓资，均系陈君独秀代为经理。现闻北大方面已议给恤金若于，而学生中亦已提议为开追悼会云。」刘文典、陈钟凡等协助料理刘师培丧事。

冬，何震由刘文典陪同，护送刘师培灵柩回扬州安葬。

1920年三十五岁

春，处理完刘师培丧事后，何震重返北京，本拟削发为尼。所以，刘富曾〈亡侄师培墓志铭〉云：「己配何氏，为余女夫扬子增生何家辂胞妹。艰难中，间关相从，武昌戎马，保其先着稿本；蚕丛崎岖，寻夫蜀道。今者嫠室哀吟，苦空彻悟，爱访名山，将为比丘终焉。」柳亚子〈南社纪略〉称何震「后来削发为尼，法名小器」。但到北京后不久，何震精神失常。柳亚子〈南社纪略〉云：「申叔死后，志剑神经病发作，曾在北大校门伏地痛哭。」本年，刘富曾、刘显曾致函陈钟凡，26云:「前以舍侄申叔疾终京寓，渥承照料一切。其后，舍侄妇神经暴发，重荷鼎力维持，俾死者得正首邱，生者得归故里。缅怀高义，腑肺铭之。」关于其精神失常

\vfill
\begin{BookNotes}
\BookNote{26.}{见吴新雷等编，《清晖山馆友声集》，江苏古籍出版社，2000年，页552。然《清晖山馆友声集》推测此函作于1919年，但刘师培棺柩1919年冬季送回扬州，何震精神病发作则在次年春回北京之后，故此函当作于1920年。}
\end{BookNotes}

% 原书第 539 页（PDF 第 552 页）
\BookPage{539}{552}
\noindent 的原因，据何氏后人何叔余先生说，刘师培逝世时，曾留下唯一一件值钱的金佛像，何震将此物存放于一位熟人家。等料理完刘师培丧事后，何震返回北京，对方不承认代为保管金佛像。因当时没有收据，所以何震气急攻心，引发精神错乱。后来，刘师培的学生帮助筹集资金，将何震送入北京的一所精神病院，其后就在精神病院逝世。何震去世的时间虽不能明，但至少应在1920年之后。


% 原书第 540 页（PDF 第 553 页）
\BookPage{540}{553}
